\documentclass[11pt]{article}

\usepackage{filecontents}
\begin{filecontents*}[overwrite]{chemical-identity-is-not-causation.bib}
@article{baronkenny1986moderator,
  author  = {Baron, Reuben M. and Kenny, David A.},
  title   = {The Moderator-Mediator Variable Distinction in Social Psychological Research: Conceptual, Strategic, and Statistical Considerations},
  journal = {Journal of Personality and Social Psychology},
  year    = {1986},
  volume  = {51},
  number  = {6},
  pages   = {1173--1182},
  doi     = {10.1037/0022-3514.51.6.1173}
}

@article{xu2026postpartum,
  author  = {Xu, Min and Luo, Yuting and Huang, Yuehua and Liu, Yunxia and Ding, Lingling},
  title   = {Risk Factors for Postpartum Depression: An Umbrella Review},
  journal = {Frontiers in Public Health},
  year    = {2026},
  volume  = {13},
  pages   = {1714668},
  doi     = {10.3389/fpubh.2025.1714668}
}

@article{ungar1968chemical,
  author  = {Ungar, George and Galvan, L. and Clark, R. H.},
  title   = {Chemical Transfer of Learned Fear},
  journal = {Nature},
  year    = {1968},
  volume  = {217},
  number  = {5135},
  pages   = {1259--1261},
  doi     = {10.1038/2171259a0}
}

@article{setlow1997ungar,
  author  = {Setlow, Barry},
  title   = {Georges Ungar and Memory Transfer},
  journal = {Journal of the History of the Neurosciences},
  year    = {1997},
  volume  = {6},
  number  = {2},
  pages   = {181--192},
  doi     = {10.1080/09647049709525701}
}

@article{gershon2013serotonin,
  author  = {Gershon, Michael D.},
  title   = {5-Hydroxytryptamine (Serotonin) in the Gastrointestinal Tract},
  journal = {Current Opinion in Endocrinology, Diabetes and Obesity},
  year    = {2013},
  volume  = {20},
  number  = {1},
  pages   = {14--21},
  doi     = {10.1097/MED.0b013e32835bc703}
}

@article{zheng2026insulin,
  author  = {Zheng, Yilin and Fang, Xiaoyi and Gao, Jiangtao},
  title   = {Re-engineering Insulin for Oral Delivery: Structural Modifications, Advanced Formulation Strategies, and Future Directions},
  journal = {Drug Delivery},
  year    = {2026},
  volume  = {33},
  number  = {1},
  doi     = {10.1080/10717544.2026.2624192}
}

@article{wiedemann2020cephalic,
  author  = {Wiedemann, Sophia J. and Rachid, Leila and Illigens, Ben and B\"oni-Schnetzler, Marianne and Donath, Marc Y.},
  title   = {Evidence for Cephalic Phase Insulin Release in Humans: A Systematic Review and Meta-Analysis},
  journal = {Appetite},
  year    = {2020},
  volume  = {155},
  pages   = {104792},
  doi     = {10.1016/j.appet.2020.104792}
}

@article{pullicin2021cephalic,
  author  = {Pullicin, Alexa J. and Glendinning, John I. and Lim, Juyun},
  title   = {Cephalic Phase Insulin Release: A Review of Its Mechanistic Basis and Variability in Humans},
  journal = {Physiology \& Behavior},
  year    = {2021},
  volume  = {239},
  pages   = {113514},
  doi     = {10.1016/j.physbeh.2021.113514}
}

@article{ader1975conditioned,
  author  = {Ader, Robert and Cohen, Nicholas},
  title   = {Behaviorally Conditioned Immunosuppression},
  journal = {Psychosomatic Medicine},
  year    = {1975},
  volume  = {37},
  number  = {4},
  pages   = {333--340},
  doi     = {10.1097/00006842-197507000-00007}
}

@article{ader1982conditioned,
  author  = {Ader, Robert and Cohen, Nicholas},
  title   = {Behaviorally Conditioned Immunosuppression and Murine Systemic Lupus Erythematosus},
  journal = {Science},
  year    = {1982},
  volume  = {215},
  number  = {4539},
  pages   = {1534--1536},
  doi     = {10.1126/science.7063864}
}

@article{vandenberg2012lowdose,
  author  = {Vandenberg, Laura N. and Colborn, Theo and Hayes, Tyrone B. and Heindel, Jerrold J. and Jacobs, David R. and Lee, Duk-Hee and Shioda, Toshi and Soto, Ana M. and vom Saal, Frederick S. and Welshons, Wade V. and Zoeller, R. Thomas and Myers, John Peterson},
  title   = {Hormones and Endocrine-Disrupting Chemicals: Low-Dose Effects and Nonmonotonic Dose Responses},
  journal = {Endocrine Reviews},
  year    = {2012},
  volume  = {33},
  number  = {3},
  pages   = {378--455},
  doi     = {10.1210/er.2011-1050}
}

@article{rhomberg2012lowdose,
  author  = {Rhomberg, Lorenz R. and Goodman, Julie E.},
  title   = {Low-Dose Effects and Nonmonotonic Dose-Responses of Endocrine Disrupting Chemicals: Has the Case Been Made?},
  journal = {Regulatory Toxicology and Pharmacology},
  year    = {2012},
  volume  = {64},
  number  = {1},
  pages   = {130--133},
  doi     = {10.1016/j.yrtph.2012.06.015}
}

@article{mar2008fiction,
  author  = {Mar, Raymond A. and Oatley, Keith},
  title   = {The Function of Fiction Is the Abstraction and Simulation of Social Experience},
  journal = {Perspectives on Psychological Science},
  year    = {2008},
  volume  = {3},
  number  = {3},
  pages   = {173--192},
  doi     = {10.1111/j.1745-6924.2008.00073.x}
}

@article{oatley2016fiction,
  author  = {Oatley, Keith},
  title   = {Fiction: Simulation of Social Worlds},
  journal = {Trends in Cognitive Sciences},
  year    = {2016},
  volume  = {20},
  number  = {8},
  pages   = {618--628},
  doi     = {10.1016/j.tics.2016.06.002}
}

@article{taylor1998imagination,
  author  = {Taylor, Shelley E. and Pham, Lien B. and Rivkin, Inna D. and Armor, David A.},
  title   = {Harnessing the Imagination: Mental Simulation, Self-Regulation, and Coping},
  journal = {American Psychologist},
  year    = {1998},
  volume  = {53},
  number  = {4},
  pages   = {429--439},
  doi     = {10.1037/0003-066X.53.4.429}
}

@article{rivkin1999mental,
  author  = {Rivkin, Inna D. and Taylor, Shelley E.},
  title   = {The Effects of Mental Simulation on Coping with Controllable Stressful Events},
  journal = {Personality and Social Psychology Bulletin},
  year    = {1999},
  volume  = {25},
  number  = {12},
  pages   = {1451--1462},
  doi     = {10.1177/01461672992510002}
}

@article{kleinman1978culture,
  author  = {Kleinman, Arthur and Eisenberg, Leon and Good, Byron},
  title   = {Culture, Illness, and Care: Clinical Lessons from Anthropologic and Cross-Cultural Research},
  journal = {Annals of Internal Medicine},
  year    = {1978},
  volume  = {88},
  number  = {2},
  pages   = {251--258},
  doi     = {10.7326/0003-4819-88-2-251}
}

@incollection{leventhal1980common,
  author    = {Leventhal, Howard and Meyer, Daniel and Nerenz, David},
  title     = {The Common Sense Representation of Illness Danger},
  booktitle = {Contributions to Medical Psychology},
  editor    = {Rachman, Stanley},
  publisher = {Pergamon Press},
  location  = {New York},
  year      = {1980},
  volume    = {2},
  pages     = {7--30}
}

@article{hagger2017commonsense,
  author  = {Hagger, Martin S. and Koch, Severine and Chatzisarantis, Nikos L. D. and Orbell, Sheina},
  title   = {The Common Sense Model of Self-Regulation: Meta-Analysis and Test of a Process Model},
  journal = {Psychological Bulletin},
  year    = {2017},
  volume  = {143},
  number  = {11},
  pages   = {1117--1154},
  doi     = {10.1037/bul0000118}
}

@article{hagger2022conceptual,
  author  = {Hagger, Martin S. and Orbell, Sheina},
  title   = {The Common Sense Model of Illness Self-Regulation: A Conceptual Review and Proposed Extended Model},
  journal = {Health Psychology Review},
  year    = {2022},
  volume  = {16},
  number  = {3},
  pages   = {347--377},
  doi     = {10.1080/17437199.2021.1878050}
}

@article{ashar2017placebo,
  author  = {Ashar, Yoni K. and Chang, Luke J. and Wager, Tor D.},
  title   = {Brain Mechanisms of the Placebo Effect: An Affective Appraisal Account},
  journal = {Annual Review of Clinical Psychology},
  year    = {2017},
  volume  = {13},
  pages   = {73--98},
  doi     = {10.1146/annurev-clinpsy-021815-093015}
}

@article{becker2023stress,
  author  = {Becker, Sara and Spinath, Birgit and Ditzen, Beate and D\"orfler, Tobias},
  title   = {Psychological Stress = Physiological Stress? An Experimental Study With Prospective Teachers},
  journal = {Journal of Psychophysiology},
  year    = {2023},
  volume  = {37},
  number  = {1},
  pages   = {12--24}
}

@article{benedetti2011placebo,
  author  = {Benedetti, Fabrizio and Amanzio, Martina and Rosato, Roberto and Blanchard, Chris},
  title   = {Nonopioid Placebo Analgesia Is Mediated by {CB1} Cannabinoid Receptors},
  journal = {Nature Medicine},
  year    = {2011},
  volume  = {17},
  pages   = {1228--1230},
  doi     = {10.1038/nm.2435}
}

@article{best2010developmental,
  author  = {Best, John R. and Miller, Patricia H.},
  title   = {A Developmental Perspective on Executive Function},
  journal = {Child Development},
  year    = {2010},
  volume  = {81},
  number  = {6},
  pages   = {1641--1660},
  doi     = {10.1111/j.1467-8624.2010.01499.x}
}

@misc{bcgov2013chicken,
  author  = {{Government of British Columbia}},
  title   = {The Real Story on How {B.C.} Chicken Is Produced},
  year    = {2013},
  url     = {https://news.gov.bc.ca/releases/2013AGRI0008-000243},
  urldate = {2026-09-12}
}

@article{campbell2012psychological,
  author  = {Campbell, J. and Ehlert, U.},
  title   = {Acute Psychosocial Stress: Does the Emotional Stress Response Correspond With Physiological Responses?},
  journal = {Psychoneuroendocrinology},
  year    = {2012},
  volume  = {37},
  number  = {8},
  pages   = {1111--1134},
  doi     = {10.1016/j.psyneuen.2011.12.010}
}

@article{danbolt2014glutamate,
  author  = {Zhou, Y. and Danbolt, N. C.},
  title   = {Glutamate as a Neurotransmitter in the Healthy Brain},
  journal = {Journal of Neural Transmission},
  year    = {2014},
  volume  = {121},
  pages   = {799--817},
  doi     = {10.1007/s00702-014-1180-8}
}

@misc{fda2012msg,
  author  = {{U.S. Food and Drug Administration}},
  title   = {Questions and Answers on Monosodium Glutamate ({MSG})},
  year    = {2012},
  url     = {https://www.fda.gov/food/food-additives-petitions/questions-and-answers-monosodium-glutamate-msg},
  urldate = {2026-09-12}
}

@article{garcia1955conditioned,
  author  = {Garcia, John and Kimeldorf, Donald J. and Koelling, Robert A.},
  title   = {Conditioned Aversion to Saccharin Resulting from Exposure to Gamma Radiation},
  journal = {Science},
  year    = {1955},
  volume  = {122},
  number  = {3160},
  pages   = {157--158},
  doi     = {10.1126/science.122.3160.157}
}

@article{garcia1966relation,
  author  = {Garcia, John and Koelling, Robert A.},
  title   = {Relation of Cue to Consequence in Avoidance Learning},
  journal = {Psychonomic Science},
  year    = {1966},
  volume  = {4},
  number  = {3},
  pages   = {123--124},
  doi     = {10.3758/BF03342209}
}

@article{zellner1983conditioned,
  author  = {Zellner, Debra A. and Rozin, Paul and Aron, Michael and Kulish, Carol},
  title   = {Conditioned Enhancement of Human's Liking for Flavor by Pairing with Sweetness},
  journal = {Learning and Motivation},
  year    = {1983},
  volume  = {14},
  number  = {3},
  pages   = {338--350},
  doi     = {10.1016/0023-9690(83)90021-8}
}

@incollection{rozin1985pavlovian,
  author    = {Rozin, Paul and Zellner, Debra},
  title     = {The Role of Pavlovian Conditioning in the Acquisition of Food Likes and Dislikes},
  booktitle = {Experimental Assessments and Clinical Applications of Conditioned Food Aversions},
  editor    = {Braveman, Norman S. and Bronstein, Pauline},
  series    = {Annals of the New York Academy of Sciences},
  volume    = {443},
  year      = {1985},
  pages     = {189--202},
  publisher = {New York Academy of Sciences}
}

@article{geha2000multicenter,
  author  = {Geha, Raif S. and Beiser, Allen and Ren, Christina and Patterson, Roger and Greenberger, Paul A. and Grammer, Leslie C. and Ditto, Anne Marie and Harris, Kathleen E. and Shaughnessy, Michael A. and Yarnold, Paul R. and Corren, Jonathan and Saxon, Andrew},
  title   = {Review of Alleged Reaction to Monosodium Glutamate and Outcome of a Multicenter Double-Blind Placebo-Controlled Study},
  journal = {Journal of Nutrition},
  year    = {2000},
  volume  = {130},
  number  = {4S Suppl},
  pages   = {1058S--1062S},
  doi     = {10.1093/jn/130.4.1058S}
}

@article{hawkins2009bbb,
  author  = {Hawkins, Richard A.},
  title   = {The Blood--Brain Barrier and Glutamate},
  journal = {American Journal of Clinical Nutrition},
  year    = {2009},
  volume  = {90},
  number  = {3},
  pages   = {867S--874S},
  doi     = {10.3945/ajcn.2009.27462BB}
}

@misc{healthcanada2022hormones,
  author  = {{Health Canada}},
  title   = {Questions and Answers: Hormonal Growth Promoters},
  year    = {2022},
  url     = {https://www.canada.ca/en/health-canada/services/drugs-health-products/veterinary-drugs/factsheets-faq/hormonal-growth-promoters.html},
  urldate = {2026-09-12}
}

@article{hoover1994sugar,
  author  = {Hoover, Daniel W. and Milich, Richard},
  title   = {Effects of Sugar Ingestion Expectancies on Mother--Child Interactions},
  journal = {Journal of Abnormal Child Psychology},
  year    = {1994},
  volume  = {22},
  number  = {4},
  pages   = {501--515},
  doi     = {10.1007/BF02168088}
}

@article{hou2017music,
  author  = {Hou, Jiancheng and Song, Bei and Chen, Andrew C. N. and Sun, Chi and Zhou, Jiaxian and Zhu, Hu and Beauchaine, Theodore P.},
  title   = {Review on Neural Correlates of Emotion Regulation and Music: Implications for Emotion Dysregulation},
  journal = {Frontiers in Psychology},
  year    = {2017},
  volume  = {8},
  pages   = {501},
  doi     = {10.3389/fpsyg.2017.00501}
}

@article{knack2024behavioral,
  author  = {Botzet, Laura Josephine and Kordsmeyer, Tobias L. and Ostermann, Sabine and Ruß, Johannes and Penke, Lars},
  title   = {Behavioural Endocrinology in the Social Sciences},
  journal = {Kölner Zeitschrift für Soziologie und Sozialpsychologie},
  year    = {2024},
  volume  = {76},
  number  = {3},
  pages   = {649--680},
  doi     = {10.1007/s11577-024-00945-3}
}

@book{marks1987book,
  author    = {Marks, Isaac M.},
  title     = {Fears, Phobias, and Rituals: Panic, Anxiety, and Their Disorders},
  publisher = {Oxford University Press},
  location  = {New York},
  year      = {1987}
}

@article{marks1987fear,
  author  = {Marks, Isaac},
  title   = {The Development of Normal Fear: A Review},
  journal = {Journal of Child Psychology and Psychiatry},
  year    = {1987},
  volume  = {28},
  number  = {5},
  pages   = {667--697},
  doi     = {10.1111/j.1469-7610.1987.tb01552.x}
}

@article{obayashi2016headache,
  author  = {Obayashi, Yoko and Nagamura, Yoichi},
  title   = {Does Monosodium Glutamate Really Cause Headache? A Systematic Review of Human Studies},
  journal = {Journal of Headache and Pain},
  year    = {2016},
  volume  = {17},
  pages   = {54},
  doi     = {10.1186/s10194-016-0639-4}
}

@article{hawkins2016glutamate,
  author  = {Hawkins, Richard A. and Vi\~na, Juan R.},
  title   = {How Glutamate Is Managed by the Blood--Brain Barrier},
  journal = {Biology},
  year    = {2016},
  volume  = {5},
  number  = {4},
  pages   = {37},
  doi     = {10.3390/biology5040037}
}

@article{roney2016frameworks,
  author  = {Roney, James R.},
  title   = {Theoretical Frameworks for Human Behavioral Endocrinology},
  journal = {Hormones and Behavior},
  year    = {2016},
  volume  = {84},
  pages   = {97--110},
  doi     = {10.1016/j.yhbeh.2016.06.004}
}

@article{smith2000transport,
  author  = {Smith, Quentin R.},
  title   = {Transport of Glutamate and Other Amino Acids at the Blood--Brain Barrier},
  journal = {Journal of Nutrition},
  year    = {2000},
  volume  = {130},
  number  = {4S Suppl},
  pages   = {1016S--1022S},
  doi     = {10.1093/jn/130.4.1016S}
}

@article{mosby2009wonton,
  author  = {Mosby, Ian},
  title   = {`That Won-Ton Soup Headache': The Chinese Restaurant Syndrome, {MSG} and the Making of American Food, 1968--1980},
  journal = {Social History of Medicine},
  year    = {2009},
  volume  = {22},
  number  = {1},
  pages   = {133--151},
  doi     = {10.1093/shm/hkn098}
}

@article{lemesurier2017uptaking,
  author  = {LeMesurier, Jennifer L.},
  title   = {Uptaking Race: Genre, {MSG}, and Chinese Dinner},
  journal = {Poroi},
  year    = {2017},
  volume  = {12},
  number  = {2},
  pages   = {1--23},
  doi     = {10.13008/2151-2957.1253}
}

@article{wolraich1995sugar,
  author  = {Wolraich, Mark L. and Wilson, David B. and White, J. Wade},
  title   = {The Effect of Sugar on Behavior or Cognition in Children: A Meta-analysis},
  journal = {JAMA},
  year    = {1995},
  volume  = {274},
  number  = {20},
  pages   = {1617--1621},
  doi     = {10.1001/jama.1995.03530200053037}
}
\end{filecontents*}


\usepackage{fontspec}
\setmainfont{NimbusRoman-Regular.otf}[
  BoldFont=NimbusRoman-Bold.otf,
  ItalicFont=NimbusRoman-Italic.otf,
  BoldItalicFont=NimbusRoman-BoldItalic.otf]
\usepackage[margin=1.15in]{geometry}
\usepackage{microtype}
\usepackage{amsmath,amssymb,amsthm}
\usepackage{booktabs}
\usepackage{enumitem}
\usepackage{titlesec}
\usepackage{csquotes}
\usepackage{xcolor}
\usepackage{fancyhdr}
\usepackage[hidelinks]{hyperref}
\usepackage[numbers,sort&compress]{natbib}

\definecolor{ink}{HTML}{222222}
\color{ink}
\titleformat{\section}{\normalfont\Large\bfseries}{\thesection}{0.75em}{}
\titleformat{\subsection}{\normalfont\large\bfseries}{\thesubsection}{0.75em}{}
\setlist{nosep,leftmargin=2em}
\emergencystretch=3em
\pagestyle{fancy}
\fancyhf{}
\fancyhead[L]{\small \textit{Chemical Identity Is Not Causation}}
\fancyhead[R]{\small Flyxion}
\fancyfoot[C]{\thepage}
\setlength{\headheight}{14pt}

\newtheorem*{principle}{Principle}

\title{\textbf{Chemical Identity Is Not Causation}\\[0.4em]
\large Molecules, Context, and the False Grammar of Bodily Explanation}
\author{Flyxion\\\small Independent Researcher}
\date{September 2026}

\begin{document}
\maketitle

\begin{abstract}
Public explanations of food, hormones, mood, and conduct often move directly from the name of a molecule to a claimed bodily or behavioral effect. Glutamate is a neurotransmitter; therefore dietary monosodium glutamate must excite the brain. A food contains hormones; therefore it must reorganize the consumer's endocrine state. Cortisol participates in stress; therefore a felt state of stress reveals what cortisol is doing. These arguments confuse chemical identity with situated causation, and often further confuse the existence of a causal contribution with its being the complete or governing explanation. A biological effect depends not only upon what a substance is, but upon dose, form, route, timing, metabolism, compartment, transport, receptor availability, regulatory capacity, developmental history, expectation, and social context, and its causal role may be that of a necessary condition, a permissive condition, a vulnerability factor, a modulator, a trigger, a mediator, or merely a correlated marker. This essay reconstructs causation as a gated sequence and a taxonomy of roles rather than a lexical association. Monosodium glutamate provides the principal chemical case; false news about a lottery ticket supplies the inverse case, in which an interpretation can produce a real bodily response without the attributed external event being real; conditioned taste aversion supplies a case where strong single-exposure inference is in fact warranted, clarifying by contrast why most single-exposure chemical attributions are not. The essay then examines the absence of counterfactuals from first-person experience, expectancy effects in judgments of children's behavior, self-confirming interpersonal loops, prospective, concurrent, and retrospective regulation, the non-monotonic development of regulatory capacities, and the difference between identifying a possible hazard and ranking an actual risk. It concludes that folk chemical explanations persist partly because they perform psychological and social work: they name diffuse suffering, externalize it, and offer a tractable intervention. A correction must therefore replace not only a false causal proposition but also the orientation that proposition supplied.
\end{abstract}
\newpage

\tableofcontents
\newpage

\section{The Molecule and the Lottery Ticket}

Consider two ordinary claims. In the first, a person eats soup seasoned with monosodium glutamate and later develops a headache. Glutamate is the major excitatory neurotransmitter in the mammalian brain, so the explanation appears almost self-evident: the soup supplied glutamate; glutamate excites neurons; the additive therefore produced the headache by disturbing brain chemistry. In the second case, a person misreads a lottery ticket, becomes convinced for several minutes that an enormous prize has been won, and undergoes an unmistakable bodily transformation. Heart rate may change; hands may shake; attention narrows; breathing changes; the imagined future reorganizes the present. The ticket is not a winner. Nothing has entered the body. Yet the response is not therefore imaginary in the sense of unreal.

The cases invert one another. The first begins with a real chemical exposure but does not thereby establish the distinctive effect attributed to it. The second begins with a false interpretation of the external world but can nevertheless produce a real bodily event. Together they separate three things that everyday explanation habitually fuses: the presence of a substance or event, the occurrence of a bodily state, and the attribution of that state to a particular cause.

It would be easy to spoil the second example by announcing that adrenaline and cortisol ``spiked.'' Unless those variables were measured, this would reproduce the mistake under examination. An intense experience does not provide introspective access to its endocrine implementation. Even laboratory stress paradigms often find imperfect correspondence between reported stress and measured cortisol response 
\citep{campbell2012psychological,becker2023stress}. The defensible claim is deliberately less specific: an appraisal can alter physiology even when the external condition being appraised is absent. Research on placebo and nocebo effects, expectancy, and appraisal supplies many measured versions of that general relation, but the lottery ticket itself remains a thought experiment rather than a hormonal assay 
\citep{benedetti2011placebo,ashar2017placebo}.

This is worth stating plainly, because placebo findings are routinely misdescribed even by people sympathetic to them: a placebo response is not evidence that ``nothing happened.'' If expectation, conditioning, clinical ritual, or prior learning measurably alters pain, nausea, autonomic activity, or treatment adherence, something happened, and it was real. What fails in a placebo response is not the effect but the attribution: the theory that a specific pharmacological ingredient produced it. The lottery ticket and the placebo response are the same shape of case. An effect can be entirely real while the causal theory attached to it is entirely false, and the reverse asymmetry holds too: a pharmacologically active treatment can work even in a patient who neither feels nor correctly identifies its mechanism, so subjective conviction and biological efficacy can come apart in either direction. Treating expectation effects as unreal shadows of ``real'' biological effects gets this backward. Expectation effects are biological events; they simply have informational and social causes rather than pharmacological ones.

This yields the principle on which the essay turns.

\begin{principle}[Causal Nonidentity]
The reality of a bodily response does not establish the reality of its attributed cause, and the presence of a substance does not establish the response attributed to it.
\end{principle}

The name is used here in a narrow, local sense and should not be confused with the unrelated ``nonidentity problem'' of population ethics or with debates over personal identity over time; the term names only the gap between chemical identity and situated causal effect that this essay is concerned with.

The principle does not deny chemistry. It denies that a chemical name is already a causal explanation. It also does not reduce bodily response to belief, as if interpretation floated free of physiology. Appraisal is itself embodied activity. The point is that neither a named molecule nor a named feeling identifies, by itself, the path connecting an exposure, an organism, and an outcome.

\section{One Error Pointed in Opposite Directions}

Two familiar attitudes appear to oppose one another. Chemical essentialism treats a manufactured, unfamiliar, or multisyllabic substance as dangerous because of what it is. Endogenous innocence treats a naturally occurring or bodily substance as safe because it is already ``part of us.'' The first supports warnings that MSG is an alien excitotoxin. The second supports the reply that glutamate is ubiquitous in food and metabolism and therefore could not matter.

These are not mirror-image doctrines so much as one error pointed in opposite directions. Both substitute identity for context. ``It is a chemical, therefore dangerous'' and ``it is natural, therefore harmless'' try to infer an effect from a noun. But water, sodium, glucose, potassium, oxygen, iron, estradiol, and glutamate are ordinary constituents of life whose consequences depend upon quantities, locations, rates, combinations, and regulatory conditions. Endogenous substances can be dangerous in the wrong compartment or concentration. Synthetic substances can be inert at relevant exposures. Familiarity is not protection, and unfamiliarity is not toxicity.

The better distinction is between \emph{substantial causation} and \emph{situated causation}. Substantial causation imagines that the power to produce an outcome inheres in an isolated substance and accompanies it everywhere. Situated causation locates an effect in a relation among material identity, exposure, pathway, target, organism, and time. This does not imply that substances have no stable properties. Without such properties, prediction would be impossible. It means that a stable chemical property is only one term in a biological event.

Language encourages the substantial error. A single word, \emph{glutamate}, names material encountered in tomatoes, aged cheese, seaweed broth, intestinal metabolism, blood, astrocytes, synaptic vesicles, and extracellular fluid near neuronal receptors. Repetition of the name creates an impression of causal continuity. Yet sameness of name cannot erase the physical separations among these sites. A map that prints the word \emph{water} over a reservoir, a pipe, a bloodstream, a cell, and the ocean has not shown that water moves freely among them.

\section{What Kind of Cause?}

Saying that a molecule or a hormone ``causes'' an outcome can be true in a weak sense while remaining profoundly misleading, because the word cause covers several distinct roles that this essay has so far treated together. Separating them lets the criticism developed above become considerably more exact. The mistake this essay is tracking is not always the claim that a causal relation exists. It is more often the promotion of one genuine causal contribution into the complete or governing explanation, at the expense of the pathway, the mediating steps, and the alternatives that were available at the moment of action.

A \emph{necessary condition} must be present for an outcome to occur at all, though its presence guarantees nothing; oxygen is necessary for a fire without being what makes one fire larger than another. A \emph{permissive condition} does not drive an outcome but must be satisfied for other causes to operate, the way adequate lighting permits reading without explaining what a person chooses to read. A \emph{vulnerability factor} raises the probability or severity of an outcome without producing it on its own, as a genetic predisposition can raise the probability of a disease that still requires some further trigger. A \emph{modulator} adjusts the magnitude or character of an already-active process without initiating it, as a hormone can sharpen or blunt an emotional response that some other event set in motion. A \emph{trigger} supplies the proximate occasion, often against a background already prepared by permissive and vulnerability factors, as a single argument can trigger a departure that months of accumulated resentment had already made likely. A \emph{mediator} lies on the causal path between an upstream cause and a downstream outcome, transmitting an effect that the upstream variable does not produce directly, a distinction with a long history in psychological methodology \citep{baronkenny1986moderator}. An \emph{amplifier} or \emph{feedback signal} feeds part of an outcome back into the system that produced it, strengthening or dampening a process already underway. A \emph{consequence} follows from the outcome rather than preceding it, and is sometimes mistaken for its cause when only their sequence is observed. A \emph{correlated marker} tracks an outcome without being causally connected to it at all, sharing a common upstream cause instead.

These roles are not interchangeable, and a claim can be true under one description and false under another. Consider a hormonal transition that alters sleep and thermoregulation; the altered sleep reduces a person's regulatory capacity the following day; a hostile encounter at work then produces an angry outburst. The hormone belongs to the causal history in a real sense, functioning here as a vulnerability factor operating through a mediator (disrupted sleep) that lowered a threshold. But the sentence ``hormones caused the outburst'' collapses this entire chain into a single term, and in doing so it suppresses the proximate trigger, the mediating pathway, the person's prior history with that colleague, and the alternative responses that remained available at the moment of the encounter. The hormone did not write the outburst; it altered a threshold that a separate, socially situated event then crossed.

This taxonomy also clarifies why removing a genuine contributing cause does not always fix a problem, and why fixing a problem does not always indict the cause that was removed. If poor sleep is the mediator through which a vulnerability factor produces a symptom, then improving sleep may resolve the symptom without the original vulnerability factor changing at all, which can create the false impression that the vulnerability factor was never a cause. Conversely, removing a trigger without addressing the vulnerability and permissive conditions that made the trigger effective may produce only temporary relief, since a different trigger will eventually find the same prepared ground. Evaluating a causal claim about MSG, a hormone, or an expectation therefore requires specifying which of these roles is being asserted, not simply whether some role exists.

\subsection{The measurable variable is not automatically the governing variable}

A closely related error deserves its own name, because it recurs wherever a quantity happens to be easy to measure. Cortisol, inflammatory markers, heart rate, and blood glucose are routinely treated as hidden agents simply because they correlate with an unwanted state, but a biomarker can occupy several different positions in a causal system: it can be upstream of the outcome, downstream as a consequence, compensatory as part of the system's own correction, incidental, or a feedback signal folding part of the outcome back into the process that produced it. Fever is the simplest illustration. An infection can cause a fever, and the fever can in turn affect experience and physiology, but the measured temperature is not identical with the underlying pathogen or with the immune response as a whole; lowering the temperature can relieve suffering without touching the infection, and in some circumstances suppressing the fever can interfere with a response that was doing protective work. Cortisol, discussed at length above, is structurally similar: it can participate in a response to a demand while also functioning as part of the negative-feedback loop that regulates that very response, so a measured cortisol profile can be as much a consequence of the conditions being called ``stress'' as their cause. This yields a principle worth stating plainly: a measurable variable is not automatically the governing variable. Quantification improves the evidence about what was measured; it does not by itself settle what causal role the measurement occupies.

The history of scotophobin is worth a brief detour here because it shows how far this particular error can travel when a research program is organized around it rather than merely tempted by it. In the late 1960s, the researcher Georges Ungar and collaborators reported that brain extracts taken from rats trained to avoid darkness could, when injected into untrained animals, transfer that dark avoidance to the new subjects \citep{ungar1968chemical}. The active component was subsequently named scotophobin, from the Greek for fear of darkness, and was discussed in some quarters as though a specific learned fear had been isolated, chemically encoded, and physically relocated from one brain to another. The claim drew substantial attention, but it did not survive as an accepted account of memory: subsequent replication attempts and methodological scrutiny undermined the original transfer effect, and a later historical review of Ungar's research program traces both its early reception and its eventual collapse \citep{setlow1997ungar}. It is worth being precise about which error scotophobin illustrates, since it is not quite the same one the fever and cortisol cases above were making. Those cases involve a real, measurable signal whose causal role, upstream, downstream, mediating, or feedback, was misidentified. Scotophobin's error came earlier and ran deeper: before the transfer effect had even been reliably established, its own name had already promoted a hypothesized molecule into the identity of a psychological state, fear of darkness encoded as a specific peptide, rather than treating it as a candidate correlate whose causal role still needed to be shown. Scotophobin is instructive not because molecules are irrelevant to memory, memory obviously depends on molecular processes at some level, but because a molecule that participates in, correlates with, or can be shown experimentally to perturb a learned behavior is not thereby identical to that behavior. ``Molecule alters or accompanies memory'' does not entail ``molecule is the memory,'' and scotophobin's very name performed a piece of illicit causal compression before the evidence had a chance to settle the question, in exactly the way ``stress hormone'' compresses an organized, multi-level response into a single measured value.

The chapters that follow return to this taxonomy repeatedly: the blood-brain barrier chapter concerns necessary and permissive conditions for glutamate to reach a target; the hormones chapter concerns the difference between a modulator and a governing cause; the expectancy chapter concerns feedback signals that a false attribution can introduce into an otherwise ordinary interaction.

\section{The Missing Middle}

The naive MSG syllogism is short:

\begin{quote}
Food contains glutamate. The brain uses glutamate as a neurotransmitter. Therefore eating glutamate directly changes glutamatergic signalling in the brain.
\end{quote}

The argument is persuasive only because it deletes nearly all of physiology. Between ingestion and behavior lie digestion, intestinal use, absorption, first-pass metabolism, systemic distribution, barrier transport, tissue access, cellular uptake, extracellular concentration, receptor subtype and location, competing signals, clearance, homeostatic response, and the relation between any local perturbation and an organism-level outcome. A more realistic causal chain is:

\begin{center}
ingestion $\rightarrow$ digestive liberation $\rightarrow$ intestinal handling $\rightarrow$ circulation\\
$\rightarrow$ tissue access $\rightarrow$ cellular access $\rightarrow$ receptor interaction\\
$\rightarrow$ physiological perturbation $\rightarrow$ experienced or behavioral outcome.
\end{center}

This is not a claim that every biological effect must follow exactly this sequence. Inhalation, injection, skin exposure, microbial metabolism, vagal signalling, immune activation, and sensory pathways take different routes. The chain instead makes visible a general requirement: the endpoint cannot be inferred merely by discovering the same molecular word near both ends.

The missing middle is also missing in a more literal sense: it omits everything that makes chemistry quantitative rather than nominal. Presence is not concentration, and a trace detectable by sensitive instruments is not the same claim as an amount capable of producing an effect. Much of what circulates is not even freely available: many molecules travel partly bound to carrier proteins, and only the unbound, free fraction is typically available to act on a receptor or cross a membrane, so a total measured level can overstate what is biologically active by a wide margin. Receptors themselves have a binding affinity and a finite number of available sites; occupying them follows its own logic of saturation, meaning that above some concentration adding more of a substance stops producing a proportionally larger effect, and below some threshold a substance may produce no detectable effect at all regardless of how reliably it is present. Two substances can also interact through antagonism, one blocking or reducing the other's effect, or synergy, together producing more than either would alone, so the presence of a substance in isolation does not fix what it will do in a body already carrying dozens of other active molecules. Kinetics adds a further dimension: how quickly a substance is absorbed, distributed, metabolized, and cleared determines how long any concentration is sustained, and a brief spike behaves differently from a sustained plateau even at the same peak value. Homeostatic and negative-feedback systems then work to return many of these values toward a set range, so a system's response to a perturbation is not simply proportional to the perturbation's size; it depends on how vigorously the system pushes back. Finally, temporal patterning matters independently of average level: many hormones are released in pulses rather than continuously, and a circadian rhythm means the same average daily concentration can correspond to very different moment-to-moment profiles depending on when in the day it was measured. None of this complexity is a reason to abandon quantitative reasoning about chemistry; it is the reason a bare claim of presence cannot substitute for it.

\subsection{Sometimes less is not merely weaker}

The dose makes the poison is a genuine improvement over folk chemical reasoning, but stated on its own it can harden into another slogan that substitutes for analysis rather than requiring it, because it quietly assumes that a dose-response relationship is monotonic: that more of a substance always produces a larger effect and less always produces a smaller one. Hormones and endocrine-active substances frequently violate this assumption. A response can increase across one range of doses and decrease across another, because receptors saturate and then become downregulated, because a negative-feedback system activates only once a threshold is crossed, because a second receptor population with an opposing effect becomes engaged only at higher concentrations, or because a compound that stimulates a process at low concentration becomes cytotoxic at high concentration through an unrelated mechanism. A substantial literature on endocrine-disrupting chemicals argues that such non-monotonic dose-response curves are common rather than rare across cell culture, animal, and human population studies \citep{vandenberg2012lowdose}, though this claim is itself contested: critics have argued that many of the cited examples amount to anecdote rather than systematic evidence, that positive and null results have not been evaluated by the same standard, and that the case for non-monotonicity as a general phenomenon remains unproven even where individual instances are plausible \citep{rhomberg2012lowdose}. The right conclusion from an active scientific dispute is not to adopt either side's confidence, but to note what both sides agree on: a dose-response shape is an empirical property of a specific system that must be measured, not a form that can be assumed in either direction.

This section functions as a deliberate stress test on the essay's own claims about dose, and it should be read that way. ``Endocrine systems can behave non-monotonically'' is not evidence that any particular substance produces any particular effect at any particular dose; it is a reason a claimed absence of effect at a tested dose cannot automatically be extrapolated to lower or higher doses without separate support, in precisely the same way a claimed presence of effect cannot be. Timing and developmental stage add a further complication that a single dose-response curve, however carefully measured, does not capture on its own: an exposure during a developmentally sensitive window is not necessarily interchangeable with the identical exposure encountered later, since the receptors, feedback systems, and target tissues an exposure interacts with are themselves still being assembled.

This discussion also disambiguates a word this essay has used in two unrelated senses. A learning or exposure sequence can be deliberately designed to increase in difficulty monotonically, each stage harder than the last, which is the sense used in the discussion of graded exposure above. A biological dose-response curve need not be monotonic at all, and assuming that it is, in either direction, is the error this section exists to name.

Each arrow in the causal chain above represents a question of admissibility. Is the material admitted across the intestinal epithelium in the proposed form and amount? Is it transformed before circulation? Does it remain available long enough to reach the relevant tissue? Can it cross the tissue interface? Does it enter the relevant extracellular or intracellular compartment? Are the necessary receptors present and accessible? Is the magnitude sufficient to escape local buffering and organism-wide regulation? Does the perturbation survive long enough, and propagate far enough, to contribute measurably to the claimed outcome?

Failure at an early gate does not logically prove that no effect of any kind occurs. A food may act through taste, expectation, gut signalling, metabolism, immune processes, or learned association even if a proposed molecule never reaches the brain. But failure does defeat the particular causal story that required passage through that gate. Biological complexity cannot be invoked as a universal rescue device for a mechanism whose required steps are absent.

\section{Compartmentalization as Admissibility}

Compartmentalization is not an inconvenience added to an otherwise simple chemistry. It is one of the means by which living chemistry becomes possible. Cells maintain gradients across membranes; organelles separate mutually incompatible processes; enzymes localize reactions; carrier proteins determine movement; binding proteins alter availability; epithelial and endothelial interfaces regulate exchange. The organism is not a bag in which all instances of a molecule meet on equal terms.

Nor does the body literally recognize provenance. A glutamate ion does not carry a label declaring that it came from MSG, Parmesan, a tomato, or protein hydrolysis. What the body can distinguish are present physical relations: position, concentration, chemical form, binding state, neighboring molecules, transport opportunity, and metabolic trajectory. Provenance matters only when its history has produced a difference in those current relations. ``Natural'' and ``artificial'' are not receptors.

The blood--brain barrier illustrates this principle sharply. It is often pictured as a wall, but a wall suggests passive exclusion. The barrier is a selective interface formed principally by specialized endothelial cells together with interacting pericytes, astrocytic endfeet, transport systems, enzymes, and regulatory signals. Different substances encounter different transport conditions. Some diffuse; some require carriers; some are actively removed; some scarcely cross under ordinary conditions; disease can alter the system.

Brain glutamate is especially revealing because the nervous system must preserve very low extracellular concentrations near cells while maintaining large intracellular and vesicular pools. Glutamate's signalling effect depends upon where it is released and which receptors it reaches, not simply upon its presence somewhere in the organism. Reviews of glutamate handling at the blood--brain barrier conclude that net entry of circulating glutamate into healthy brain tissue is highly restricted, while transport and uptake mechanisms help maintain the relevant gradients 
\citep{smith2000transport,hawkins2009bbb,danbolt2014glutamate,hawkins2016glutamate}. The brain also synthesizes and recycles its own transmitter pools. Saying ``glutamate is excitatory'' without specifying extracellular access near appropriate receptors is like saying electricity is dangerous without distinguishing charge stored in a battery from current passing through the heart.

This makes the barrier a form of biological admissibility gating. Presence in blood does not license presence in brain extracellular fluid. Presence in an endothelial cell does not license net delivery into neural tissue. Presence in tissue does not license synaptic action. Even receptor binding does not license a particular mood, judgment, or behavior. At every scale, a possible continuation must satisfy local conditions.

Admissibility is not absolute. Barrier integrity can change; circumventricular organs differ from most brain tissue; dose can overwhelm a regulatory process; development, disease, medication, and nutritional state can alter transport. The correct conclusion is therefore not ``dietary glutamate can never matter.'' It is that any claim about how it matters must specify and support the relevant path. The barrier replaces certainty in both directions with a tractable set of mechanistic questions.

\subsection{Route can reverse the apparent meaning of a substance}

The blood-brain barrier is the most vivid gate in this essay, but treating it as a special case specific to the brain would understate the argument. A substance's route of entry can change not merely the degree but the entire character of its effect: a compound can be inert when eaten and dangerous when inhaled, harmless on intact skin and consequential when injected, useful at one anatomical site and destructive at another. Insulin is the clearest familiar illustration. Ordinary insulin is a peptide, and peptides are precisely the class of molecule the gastrointestinal tract is built to dismantle: gastric acid denatures it and digestive enzymes including pepsin, trypsin, and chymotrypsin degrade it before meaningful absorption can occur, which is why decades of pharmaceutical effort have gone into engineering oral insulin formulations that survive the gut long enough to work, with unmodified oral insulin's bioavailability generally remaining below one percent \citep{zheng2026insulin}. ``Insulin is present in the body'' therefore tells us almost nothing about what swallowing an insulin-containing substance would do, because the digestive route was never a candidate pathway for that molecule to begin with.

This generalizes the barrier argument rather than restricting it to neuroscience: the skin, lungs, gut, liver, placenta, kidneys, and every cell membrane impose their own admissibility conditions, each different from the blood-brain barrier's specific transport logic. It also supplies a direct answer to a common evidential move: citing a study in which a compound was injected into an animal as evidence about the danger of eating that same compound. An injection deliberately bypasses digestion, first-pass liver metabolism, and intestinal barrier transport, all of which a dietary exposure would have to pass through. Such a study can be excellent evidence about what a compound does once it is systemically available, and correspondingly weak evidence about what will happen when the same compound is eaten, because the whole point of the missing middle developed above is that getting to systemic availability is most of what the causal question depends on.

\section{MSG as a Case Study, Not a Cause C\'el\`ebre}

Monosodium glutamate is the sodium salt of glutamic acid. In aqueous conditions it dissociates into sodium and glutamate. The glutamate is not a sinister analogue that merely resembles dietary glutamate; it is the same chemical species encountered in many foods and in ordinary metabolism. The United States Food and Drug Administration notes both that glutamate occurs naturally in foods such as tomatoes and cheeses and that the body handles glutamate from MSG in the same way as glutamate naturally present in food 
\citep{fda2012msg}. This fact defeats claims that depend upon MSG glutamate possessing a unique artificial identity. It does not, by itself, prove that every dose in every person under every condition is without effect.

The evidence is more interesting than either popular pole allows. Blinded challenge studies have produced mixed results, especially when large bolus doses were consumed without food. In one multicenter double-blind study of people who considered themselves sensitive to MSG, initial responses occurred in both MSG and placebo conditions; higher symptom scores appeared under some MSG challenges, but reproducibility across rechallenges was limited and responses were associated with large doses without food 
\citep{geha2000multicenter}. That study's proceedings appeared as a symposium supplement to \emph{The Journal of Nutrition} whose financial support came from the International Glutamate Technical Committee, an industry association \citep{geha2000multicenter}. This does not invalidate the experiment's design or its double-blind, placebo-controlled methodology, but it belongs in the evidential chain rather than being left implicit: a favorable-looking result should be traced to its funding context as routinely as an alarming one is traced to its dose and blinding. A systematic review focused on headache found inconsistent results and substantial methodological problems, including inadequate blinding caused by MSG's taste and doses or concentrations unlike ordinary culinary exposure 
\citep{obayashi2016headache}. The FDA reports that some sensitive individuals may experience temporary, generally mild symptoms after consuming three grams or more without food, while ordinary reports have not established the broad syndrome often attributed to MSG 
\citep{fda2012msg}.

The appropriate conclusion is neither ceremonial reassurance nor generalized alarm. It is disaggregation. A claim about a large aqueous bolus consumed while fasting is not automatically a claim about a fraction of that quantity dispersed through a meal. A report of transient symptoms is not evidence of neuronal excitotoxicity. A possible headache trigger for some individuals is not a general theory of behavioral disturbance. A regulatory judgment about typical use is not a promise that no individual can ever experience an adverse reaction. The dose, vehicle, meal context, endpoint, population, and study design are parts of the claim rather than qualifications appended afterward.

MSG is valuable here precisely because the essay is not ultimately about MSG. It demonstrates how a molecule can become a cultural object. Its name gathers associations with processing, foreignness, industrial food, restaurants, headache, and loss of control. The history of ``Chinese restaurant syndrome,'' named after a 1968 letter describing symptoms following a meal at a Chinese restaurant, has been examined by historians of medicine as substantially organized by a racialized discourse that framed Chinese cooking as exotic and suspect well before MSG's mechanism had been seriously investigated \citep{mosby2009wonton}. Rhetorical analysis of the original medical correspondence has traced how satire built on stereotypes of Chinese identity, offered as insider humor among physicians, was subsequently taken up in the popular press as unquestioned medical fact, with the humor dropped and the prejudicial framing retained \citep{lemesurier2017uptaking}. Free glutamate is, meanwhile, widely distributed across cuisines, present in comparable or greater concentrations in tomatoes, aged cheeses, and cured meats common to cooking traditions that were never subjected to the same suspicion. The cultural story affects which meals are monitored, which sensations become salient, and which ingredient receives blame. Chemistry and expectation then enter the same event through different routes, and neither the historical record nor the biochemistry alone accounts for why one cuisine became the syndrome's namesake.

\section{Most Serotonin Is Not Brain Serotonin}

Serotonin offers a cleaner compartmentalization case than glutamate, because public vocabulary calls it the ``happiness chemical'' while treating every serotonin-related claim, gut or brain, dietary or pharmaceutical, as though it drew on one shared reservoir. It does not. The substantial majority of the body's serotonin, on the order of ninety percent or more by most estimates, is produced outside the central nervous system, chiefly by enterochromaffin cells in the intestinal lining, where it regulates peristalsis, secretion, and local inflammatory signaling rather than mood \citep{gershon2013serotonin}. Peripheral, gut-derived serotonin does not simply cross into the brain to join central serotonergic signaling; the central and peripheral pools are produced by different enzymes, held in different compartments, and serve substantially different local functions, joined by the same molecular name without being causally continuous \citep{gershon2013serotonin}.

This should not be read as a claim that the gut cannot influence mood, and it is worth being explicit about why not, since the essay's own admissibility framework would be misapplied if it were used to sever every gut-brain connection. Communication between the two does occur, but through specific, traceable routes, vagal signaling, immune and inflammatory mediators, and metabolic byproducts among them, not through peripheral serotonin itself diffusing upward into central circuits. If a claim asserts that something happening in the gut is affecting mood or cognition, the admissibility standard developed throughout this essay applies without modification: the claim needs to specify which of these routes is doing the work, rather than being satisfied by noting that the same molecular name, serotonin, appears at both ends.

The formulation this section adds to the essay's vocabulary is close to, but usefully distinct from, the glutamate case: MSG showed that eating a neurotransmitter does not pour it directly into synapses; serotonin shows that producing a neurotransmitter somewhere in the body does not make it available everywhere its name subsequently appears in a sentence. The same molecule can belong to genuinely different causal systems, because what determines its effect is not the name but the access, source, and organization surrounding each of its separate pools.

\section{Hormones Are Conditions, Not Commands}

Hormonal explanation often repeats the MSG syllogism at the level of conduct. A hormone participates in sexual desire, stress, attachment, appetite, or reproductive physiology. A person experiences desire, fear, hunger, or a reproductive decision. The hormone is then treated as the hidden author of the action. ``Her hormones made her do it'' compresses a distributed process into a chemical command.

Hormones unquestionably matter. They can change energy, sleep, pain, thermoregulation, libido, attention, reward sensitivity, stress responsiveness, and the salience of environmental cues. Puberty, pregnancy, postpartum transitions, menstrual cycling, perimenopause, endocrine disorders, and hormone treatments can produce profound experiences. Rejecting hormonal determinism must not become a polite way of denying embodied difficulty or returning every problem to attitude.

But influence is not instruction. Hormones act through receptors distributed unevenly across tissues; effects can depend on concentration, timing, pulsatility, binding proteins, receptor density, prior exposure, interactions with other signals, and the state of neural and social systems. Behavior can also alter hormonal state. Human behavioral endocrinology increasingly treats hormone-behavior relations as reciprocal and context-sensitive rather than as one-way outputs 
\citep{roney2016frameworks,knack2024behavioral}. A hormone may alter the difficulty landscape without selecting the route through it.

\subsection{Cortisol: Which Measure, Which Moment?}

Cortisol deserves separate treatment because it is probably the single most invoked hormone in folk explanation, and the invocation almost never specifies which of several distinct things ``cortisol'' could mean. Cortisol has a basal level that follows a circadian rhythm, typically peaking soon after waking and declining across the day; it also has an acute reactivity profile, a rise and fall following a specific stressor; and it has a chronic profile, an average level or pattern sustained over weeks or months that differs from any single day's reading. ``My cortisol is high'' could refer to any of these, or to none of them in a measured sense at all, since a review focused explicitly on four steroid hormones including cortisol notes that its behavioral and psychological correlates depend on which of these temporal features is at issue \citep{knack2024behavioral}.

The relationship between subjectively felt stress and measured cortisol is, further, weaker and more variable than the folk picture assumes. A review of the psychosocial stress literature found that the correspondence between emotional stress responses and physiological stress responses, cortisol among them, is inconsistent across studies, sometimes strong and sometimes essentially absent \citep{campbell2012psychological}. An experimental study with prospective teachers illustrates the point directly: an induced stressor produced a significant rise in salivary cortisol between the stress and control groups, but the same study found no corresponding difference in self-reported psychological stress, and the two measures did not correlate within individuals \citep{becker2023stress}. A person can therefore feel highly stressed with unremarkable cortisol, or show a measurable cortisol response while reporting no distress at all. Saying ``my cortisol is high'' from how one feels is not merely unmeasured; the feeling and the hormone have been shown, experimentally, to be capable of moving independently of one another.

None of this means cortisol is causally idle. It plausibly participates in sleep disruption, appetite change, and altered threat sensitivity in ways that feed into the mediating pathways discussed in the causal-role taxonomy above. The point is narrower and more useful: a claim invoking cortisol should specify which temporal feature is meant, by what measure it was assessed, and whether the claim concerns the hormone's basal, reactive, or chronic profile, rather than treating ``cortisol'' as a single dial that a felt state can be read off directly.

This distinction matters acutely when explaining women's conduct. To say that women ``lose judgment'' under hormonal influence takes a real domain of bodily change and converts it into a license for biological fatalism. The behavioral-endocrinology literature already introduced above argues explicitly against collapsing hormonal variation into a single fixed hormone-behavior link, treating hormones instead as one input among several reciprocal, context-dependent systems \citep{roney2016frameworks,knack2024behavioral}. Falling in love, wanting to create a family with a partner, fearing the loss of fertility, remaining in a coercive relationship, declining parenthood, or revising a life plan cannot be rendered intelligible by an endocrine variable alone. These actions occur inside narratives about adulthood, obligation, identity, security, time, and imagined futures. Education, economic dependence, housing, access to care, contraception, family expectations, partner behavior, prior trauma, and available cognitive strategies can exert more explanatory force than the chemical variable used to summarize the scene.

Postpartum depression supplies a particularly important warning, and applying the causal-role taxonomy introduced earlier makes clear why. An umbrella review synthesizing dozens of meta-analyses found that the strongest identified risk factors were not exclusively or even primarily hormonal: intimate partner violence, poor social support, sleep disorders, unintended pregnancy, and a personal history of depression all carried substantial, independently estimated associations with postpartum depression, while protective factors included postpartum physical activity, doula support, and breastfeeding support \citep{xu2026postpartum}. This is not evidence that hormones are irrelevant. Rapid postpartum hormonal withdrawal plausibly functions as a vulnerability factor, and disrupted sleep plausibly functions as a mediator through which that vulnerability translates into symptoms severe enough to notice. But a person who reports distress and is refused a hearing by a partner is not experiencing a private chemical event alone: the partner's response is not a bystander to the causal picture but an active input, changing isolation, help-seeking, sleep, burden, and, per the review's own strongest single risk factor, danger. Treating a hormonal transition as the whole story forecloses attention to the risk factors the evidence actually ranks highest, several of which are remediable in ways a hormone is not. Biology and social organization are not rival explanations placed on opposite sides of the person. They are interacting parts of one causal system, with the hormonal contribution better described as a vulnerability factor operating alongside mediators, triggers, and permissive conditions than as a governing cause in its own right.

Perimenopause reveals a second complication. Long-neglected experiences can finally receive medical recognition while simultaneously becoming a commercial frontier. Better recognition can relieve suffering and correct historical dismissal. The same attention can also generate compulsory optimization: protein targets, heavy lifting, cortisol control, supplements, body-shape warnings, diagnostic content, and an expanding market of solutions. The fact that a market exaggerates a problem does not show that the problem is unreal. The fact that a transition is biological does not validate every product sold in its name. Again, identity does not decide the case.

Claims about ``hormones in food'' require equally careful separation, and ordinary conversation typically collapses three distinct propositions that need to be evaluated independently. The first is that an animal naturally contains hormones; this is true of essentially every animal product, since any animal with an endocrine system produces and circulates its own hormones as a basic condition of being alive. The second is that a producer may have administered additional hormonal growth promoters; this varies sharply by animal, jurisdiction, and specific compound; hormonal growth promoters are approved in Canada for beef cattle but not for poultry, and recombinant bovine somatotropin is not approved for dairy cattle at all \citep{healthcanada2022hormones,bcgov2013chicken}. The third, and the one that actually bears on a health claim, is that consuming the resulting product exposes a human target tissue to a biologically consequential dose; this does not follow automatically from either of the first two, since it further requires the hormone to survive digestion, retain activity, reach an absorbed and bioavailable form, and do so at a level that is large relative to the body's own endogenous production and turnover.

Chicken is a useful illustration of how this collapse creates false differentiation: a label reading ``hormone-free chicken'' can imply the animal received no artificial hormones, which in a market where such use was never permitted for poultry in the first place tells the buyer nothing that was not already true of every other chicken sold there. Milk presents the opposite complication: it is manufactured by a mammal's endocrine-regulated body and naturally contains biologically active hormone-related compounds regardless of any agricultural policy, so the relevant questions are not whether hormones are present but their chemical stability under digestion, their bioavailability, their concentration relative to what the body already produces on its own, and whether they reach a target tissue in a form and amount capable of doing anything. A jurisdiction's policy on recombinant bovine somatotropin governs the second proposition; it is not, by itself, evidence about the third. None of this is intended to reach a blanket verdict on milk or chicken specifically. Its purpose is to show that a single phrase, ``hormones in food,'' can refer to several materially different claims, and that settling one does nothing to settle the others. Presence must be followed by amount, chemical stability, digestion, bioavailability, exposure relative to endogenous production, target access, and measured outcome. ``Contains hormones'' is the beginning of a question, not its answer.

\subsection{Detection is not exposure}

Instruments now detect quantities of a substance that would have been undetectable a generation ago, and public discourse routinely converts that analytical achievement directly into a toxicological claim: a compound was detected, therefore consumers were meaningfully exposed; exposure occurred, therefore some biological perturbation followed; perturbation occurred, therefore risk increased. Each arrow in that chain is a substantive claim requiring its own support, not a restatement of the one before it. Improving the sensitivity of a detection method can make a genuinely unchanged situation look newly alarming, because what changed was the instrument's capacity to find a trace, not the amount of the substance actually present. ``Zero'' in a food-safety report typically means below that method's detection or quantification limit, not the metaphysical absence of a molecule; nearly every biologically active substance discussed in this essay, including glutamate and the endogenous hormones just discussed, is present in food and in the body at some nonzero level essentially all the time. Detecting an estrogenic compound in milk, to take the example directly at hand, is not yet evidence of a physiologically consequential exposure; the argument still needs concentration, frequency and duration of consumption, chemical stability through digestion, absorption, comparison against endogenous production, and a demonstrated target effect, all of which were already required by the disaggregation urged above. The reverse error is equally available and equally worth naming: a low detected concentration is not automatically irrelevant either, when exposure is persistent, the target tissue is unusually sensitive, or a developmental window creates a vulnerability that an adult exposure would not. Detection opens a causal inquiry. It does not, by itself, close one.

\section{The Multiplicity Hidden by a Single Name}

Folk biochemical explanation often selects one famous variable from a densely interacting system. Glutamate becomes excitation, serotonin becomes happiness, dopamine becomes pleasure, oxytocin becomes love, and cortisol becomes stress. These labels are memorable because they reduce a many-dimensional system to a moral cast of chemical characters.

The nervous and endocrine systems do not respect that cast. Neurotransmitters and neuromodulators act through multiple receptor families and subtypes with different distributions and effects. The same molecule can participate in different processes depending upon receptor, circuit, timing, and baseline state. Hormones interact with neural, immune, metabolic, and autonomic processes. Sleep, illness, recent food, circadian phase, menstrual or reproductive state, medication, learned expectation, and social setting can change together. A person who says ``my cortisol is high'' may be naming a recognizable experience, but the experience does not reveal a concentration, its temporal profile, its source, or its causal priority.

Two plain-language concepts from biology name why the one-to-one dictionary fails structurally rather than just empirically. \emph{Pleiotropy} is the property by which one signal contributes to many different outcomes: the same molecule of oxytocin participates in uterine contraction, milk letdown, social bonding, and trust-related behavior, depending on where and when it acts, which means no single behavioral gloss can be its meaning. \emph{Degeneracy} is the complementary property by which different mechanisms can produce a similar functional outcome: a feeling of calm can be produced through several distinct neurochemical routes, several distinct cognitive strategies, or several distinct social circumstances, none of which is uniquely ``the'' pathway that calm runs through. Pleiotropy defeats the idea that a molecule has one meaning; degeneracy defeats the idea that an outcome has one cause. Together they are the structural reason dopamine cannot simply mean pleasure, oxytocin cannot simply mean love, and cortisol cannot simply mean stress: each molecule does too many things, and each outcome can be reached too many ways, for the one-word gloss to survive contact with the underlying system.

This is not an argument that subjective knowledge is worthless. A person may accurately report pain, fatigue, dread, agitation, cognitive difficulty, or a recurring temporal pattern. Longitudinal familiarity with one's own body can reveal regularities missed during a short clinical visit. Testimony is evidence of experience and can be evidence about timing and association. It becomes unreliable when promoted, without an additional method, into a selective assay of hidden chemistry.

The epistemic limitation is structural. Consciousness presents a highly processed condition, not an inventory of every signal that contributed to it. Even researchers with endocrine training do not identify their circulating hormone levels by introspection. Technical vocabulary does not create access. It can instead give an underdetermined inference the appearance of measurement.

\section{A Lived Sequence Has No Control Condition}

Suppose someone says, with complete sincerity, ``I ate MSG and then developed a headache.'' The statement may accurately report a sequence. It does not contain the comparison needed to isolate the cause. The same person cannot rerun the afternoon with sleep, hydration, expectations, stress, lighting, meal composition, illness, caffeine, posture, and every other relevant variable held fixed while MSG alone is removed.

\begin{principle}[The Missing Comparison]
A lived sequence supplies temporal order, not its own comparison condition.
\end{principle}

This is why an unaided lived sequence does not supply its own causal identification, not merely because people are inattentive or biased. A single life follows one realized branch. The counterfactual branch, the same relevant world without the proposed cause, is unavailable. Causal methods are technologies for approximating that missing branch. This does not make every anecdote worthless. An anecdote can discover a hypothesis, motivate a reasonable precaution, or supply genuine evidence within a mechanism that independent research has already established; what it generally cannot do, alone, is estimate a prevalence, isolate a previously unsupported mechanism, or justify a population-level claim. Its evidential weight depends on repetition, specificity, the competing explanations it must rule out, prior mechanistic knowledge, and how distinctively the observed pattern matches a mechanism already known to produce it. The conditioned-taste-aversion case discussed later in this essay is the clearest illustration: it is a single-trial inference that succeeds evidentially, not by exception to this standard, but because it satisfies exactly these conditions.

Randomized controlled experiments distribute known and unknown influences probabilistically across groups. Blinding prevents expectations from systematically traveling with treatment assignment. Placebos approximate aspects of the treatment context. Crossover and carefully designed within-person studies compare individuals against themselves. Withdrawal and rechallenge test recurrence. Dose--response analysis asks whether the effect changes coherently with exposure. Natural experiments and instrumental-variable methods exploit structured variation that is not experimentally assigned. Mechanistic interventions test links within a proposed path. None recreates the literally identical unrealized afternoon. Each constructs a disciplined substitute with characteristic assumptions and failure modes.

Meta-analysis occupies a different position. It does not manufacture a counterfactual. It aggregates studies that attempted to do so, potentially increasing precision and revealing heterogeneity. It can also aggregate confounding, publication bias, weak blinding, incompatible endpoints, or irrelevant exposure conditions. A pyramid that automatically places meta-analysis above all individual studies mistakes a publication form for an evidential guarantee.

The absence of a private control condition does not require paralysis. People must make decisions before ideal evidence exists. A repeated, tightly timed reaction may rationally justify personal avoidance even when it cannot establish a general mechanism. The key is calibration: ``I choose not to eat this because I repeatedly feel worse afterward'' is narrower than ``this additive disrupts everyone's neurotransmitters.'' Personal policy can be prudent under uncertainty without being redescribed as universal biochemistry.

\section{Expectation Leaves the First Person}

The problem is not confined to people interpreting their own bodies. A causal category can organize the perception of another person. Hoover and Milich's 1994 experiment recruited thirty-five boys aged five to seven whose mothers believed them to be sugar-sensitive. All boys received a drink sweetened with aspartame. Mothers were told either that their child had received a large dose of sugar or that he had received placebo. Mothers in the sugar-expectancy condition rated their sons as significantly more hyperactive. Observations also found differences in maternal behavior: those mothers maintained greater physical proximity and criticized, looked at, and talked to their sons more often 
\citep{hoover1994sugar}.

The study is sometimes summarized as proof that sugar does not affect children's behavior. That overstates what its particular manipulation established, because no child received sugar. Its direct result concerns expectancy: assigned belief about sugar altered maternal ratings and interactions despite identical drinks. Broader controlled literature and a subsequent meta-analysis, published the following year, addressed the separate question of sugar's behavioral effects more directly 
\citep{wolraich1995sugar}.

The distinction matters because the expectancy study shows more than a mistaken report. The belief entered the interaction. A mother who expected chemically caused hyperactivity watched differently and acted differently. The causal story changed the surveillance applied to the child and the conduct directed toward him.

This supports, but does not fully demonstrate, a longer recursive mechanism:

\begin{center}
causal label $\rightarrow$ intensified surveillance $\rightarrow$ altered interaction\\
$\rightarrow$ interpersonal friction or changed conduct $\rightarrow$ apparent confirmation.
\end{center}

Hoover and Milich demonstrated the early portion: label, rating, and altered maternal behavior. Establishing that the altered interaction subsequently produced more difficult child behavior, which was then read back as confirmation, would require further measurement. The loop is therefore a theoretically supported extension rather than a result to be attributed wholesale to that experiment.

Its generality is nonetheless apparent. ``You seem hormonal today'' changes the interaction in which the judgment is supposedly being tested. ``You are not yourself; did you eat something?'' instructs both parties about which sensations to notice and which explanations are admissible. A clinician's warning can amplify symptom monitoring; a partner's attribution can become dismissive; a parent's expectation can become controlling. The chemical category is no longer merely a belief about an antecedent. It becomes an active component of the present social cause.

\section{Behavior Is Organized, Not Secreted}

Chemical determinism pictures action as the visible secretion of an internal state. Hunger produces eating, fatigue produces sleep, hormones produce desire, cortisol produces poor judgment. Yet people routinely reorganize the conditions under which bodily states will be encountered. This does not free behavior from biology, and the point of this section is not that behavior is unimplemented by the body but that it is not merely read off a chemical state as a direct behavioral output. It demonstrates that biological creatures can act upon models, habits, rules, environments, and anticipated costs.

Three forms of regulation should be distinguished. \emph{Prospective regulation} acts upon a forecasted state while intervention remains comparatively cheap. A person goes to bed before becoming overwhelmingly tired, eats before desperate hunger impairs choice, stops eating ice cream before fullness becomes discomfort, declines a disturbing film before sleep, delays opening a long-awaited package until obligations are complete, or avoids taking small children to frightening media immediately before bedtime. The current felt state does not command the action. A represented later state helps select it.

Prospective regulation provides a concrete account of executive function without turning it into a mysterious quantity of willpower. It is the practiced capacity to let a modeled future condition alter a present choice. The strategy may involve inhibition, working memory, attention shifting, environmental design, language, and habits acquired through instruction and repetition. Often the most effective executive act is not heroic resistance at the point of maximum temptation but arranging the environment earlier, when the decision is inexpensive.

\emph{Concurrent regulation} operates through feedback on a state already underway. Fidgeting, pacing, breathing, rhythmic movement, sensory modulation, and music may alter attention, arousal, or emotional experience in real time. They need not depend centrally upon a detailed forecast. They are control processes coupled to current feedback. Evidence concerning music and emotion regulation is suggestive but heterogeneous; music can evoke as well as regulate emotion, and neural overlap does not itself prove a therapeutic mechanism 
\citep{hou2017music}. Claims that empathy or musical experience is literally internal muscular simulation should therefore be presented as contested embodied hypotheses rather than settled facts.

\emph{Retrospective model revision} operates after an event rather than before or during it. A person updates expectations, constructs a narrative about what happened, rehearses a different response for next time, or changes an environment so a future recurrence will be handled differently. This category connects regulation directly to the causal-attribution errors examined throughout this essay, because the explanation a person settles on after an event becomes an input to which future intervention gets attached. Someone who attributes a headache to MSG may remove MSG from future meals while leaving sleep debt, meal timing, dehydration, or social anxiety, whatever the actual mediator was, untouched. The retrospective narrative is not idle commentary on a finished episode; it is the mechanism by which a mistaken causal category gets built into a person's future regulatory strategy, potentially for years, without the strategy ever being tested against the possibility that the named cause was wrong. A retrospective narrative can be revised as new information arrives, which is what separates useful reflection from a fixed diagnosis; treating a first attribution as final forecloses that revision before it has a chance to occur.

The distinction between these three forms prevents ``modeling'' from expanding until it names every adaptive act. Prospective control depends upon representing what has not yet occurred. Concurrent control depends upon sensing and modifying what is occurring. Retrospective revision depends upon reconstructing what already occurred in a way that shapes what happens next. All three can become habitual, scaffolded, and socially taught. All three can fail under pain, sleep deprivation, illness, fear, intoxication, grief, or overload.

Self-doubt belongs here only in a bounded form. Generalized self-distrust can disable agency. Prospective skepticism is narrower: ``Given this recognizable situation, I may later regret acting now, and postponement is cheap.'' It attaches doubt to a specified prediction, a known vulnerability, and a feasible precaution. Cynicism toward the self is not the virtue. Calibrated distrust of a particular present impulse can be a learned regulatory instrument.

\section{Physiology Can Act on Information Before the Exposure Arrives}

Prospective regulation, as described above, is a capacity people exercise deliberately: choosing to sleep, eat, or decline stimulation before a forecasted state arrives. It is worth showing that physiology itself can behave prospectively, at a level below anything a person decides, because this closes a gap the earlier discussion of appraisal and interpretation left open. If the body can act on evidence that an exposure is coming before the exposure itself becomes systemically available, then the lottery ticket's lesson, that an interpretation can produce a real bodily response without the corresponding external event, is not confined to consciously held beliefs. It reaches into physiological preparation that a person never explicitly forms an opinion about at all.

\subsection{The body responds before the food arrives}

The sight, smell, taste, and even the mere expectation of food can trigger physiological preparation, cephalic-phase responses, before any nutrient has been absorbed into circulation. These can include salivation, gastric secretion, pancreatic activity, and, in some studies, an early and transient rise in insulin. The human evidence on cephalic-phase insulin release is genuinely mixed rather than settled: a systematic review and meta-analysis found that the majority of studies detected some hormonal change in response to food-associated cues, but also that the effect size is small relative to insulin's ordinary spontaneous fluctuation and highly sensitive to the type of stimulus and how it was presented \citep{wiedemann2020cephalic}, and a subsequent mechanistic review emphasized that this variability, rather than any single robust number, is itself the field's central finding \citep{pullicin2021cephalic}. This essay's own standard applies to this claim as much as to any other: the section is not asserting a universal reflex, only that sensory and anticipatory information can measurably influence a hormonal variable before the supposed chemical trigger, absorbed food, has systemically arrived at all.

The point this secures is a further instance of the Principle of Causal Nonidentity, extended forward in time rather than sideways into appraisal. If an insulin change follows the sight or expectation of a sweet food, that physiological event cannot be attributed solely to absorbed glucose, because no glucose has yet been absorbed. The organism is responding to information correlated with an anticipated exposure, not to the exposure itself. This is not a mystical claim about mind over matter; cephalic-phase responses are neurally mediated, measurable, and have an evident adaptive rationale, since preparing digestive and metabolic systems in advance of a meal helps limit the disruption that ingestion would otherwise cause. It does mean that ``insulin rose because of the sugar in the food'' cannot be assumed as the default explanation for an early insulin rise without checking whether the food had actually been absorbed yet.

\subsection{When a flavor enters the immune system}

A more striking demonstration removes conscious appraisal from the picture entirely. In a now-classic set of experiments, rats were given a distinctively flavored saccharin solution paired with cyclophosphamide, a drug that suppresses immune function and also happens to produce gastrointestinal illness; predictably, the rats developed a conditioned aversion to the saccharin flavor, the ordinary taste-aversion learning discussed earlier in this essay \citep{ader1975conditioned}. What was not predicted, and what makes the result worth including here, is that later presentation of the saccharin flavor alone, with no additional cyclophosphamide, reproduced measurable immunosuppression: animals conditioned in this way and then re-exposed to the flavor showed a reduced antibody response to a subsequent antigen challenge, relative to unconditioned animals or those never re-exposed to the flavor \citep{ader1975conditioned}. A later study extended the phenomenon to mice genetically prone to a lupus-like autoimmune disease, in which the conditioned flavor cue, administered without the full dose of drug it had previously been paired with, measurably slowed disease progression and extended survival \citep{ader1982conditioned}.

This case is a stronger instance of the general point than the lottery ticket or cephalic-phase insulin, because it requires no verbal belief, no conscious expectation, and no appraisal of a situation as threatening or promising. A learned sensory cue acquired direct access to an immune regulatory pathway through ordinary associative conditioning. This dissolves any remaining temptation to treat ``psychological'' and ``biological'' as two separate substances that occasionally interact; here, learning simply is part of how the immune system is organized, in the same sense that a habit is part of how a nervous system is organized. It also extends the lottery ticket's inversion one step further. Appraisal is one route by which information becomes bodily causation. Conditioning, association, and physiological anticipation are others, and they can produce real, measurable, and in the lupus case clinically consequential effects without a person forming any explicit proposition about what is happening to them at all.

\section{The Same Dose Is Not the Same Exposure Twice}

Tolerance and withdrawal introduce a further variable this essay's causal chain has so far left implicit: history. A dose that produces a large effect on first exposure may produce a much smaller effect after repeated exposure, as receptors adapt, as compensatory systems engage, and as the organism reorganizes around the substance's expected presence. Removing the substance afterward can then produce a new effect, sometimes the opposite of the original one, not because the chemistry of the dose has changed but because the receiving system has. Caffeine is the most familiar illustration available to nearly every reader: a first exposure, a habitual daily exposure, an interrupted habitual exposure, and a resumed exposure are causally distinct events despite involving the identical named molecule at what may be the identical dose. Baseline arousal, receptor adaptation, learned timing cues, accumulated sleep debt, and the specific discomfort of withdrawal all change what ``a cup of coffee affected me'' actually refers to on a given morning.

This sharpens situated causation, developed early in this essay, by adding a temporal dimension that the earlier formulation left mostly implicit. An exposure is not an encounter between a fixed molecule and a fixed body; it is an event happening to a body that has already been altered by its previous encounters with that same molecule, or with molecules acting on overlapping systems. The same reasoning applies well beyond caffeine, to medications, alcohol, nicotine, learned food preferences, chronic stress, and habitual exercise: in each case, the relevant question is not simply what dose was received, but what history of prior exposure that dose is landing on. This is also where chemistry connects most directly to development. What an organism can do with a given exposure changes across its lifespan and its individual history, which is precisely the subject the next section takes up in more detail.

It is tempting to describe regulatory development as monotonic improvement. The description is wrong if monotonic means that every later point equals or exceeds every earlier point. Capacities develop unevenly, can plateau, depend upon task demands, and may be temporarily inaccessible. Different components of executive function follow different developmental trajectories from early childhood through adolescence 
\citep{best2010developmental}. Apparent performance can also change nonlinearly as the task itself recruits newly automatic skills.

Isaac Marks used the expression \emph{ontogenetic parade} for the recognizable emergence, prominence, and decline of normal childhood fears. Particular fears rise and fall in a developmental sequence shaped by maturation and experience 
\citep{marks1987fear,marks1987book}. The metaphor transfers productively to regulation, provided it is marked as an extension rather than Marks's own term.

Regulatory strategies have an ontogenetic succession. Immediate discharge and external soothing precede complex deliberate inhibition. Movement, repetition, and attention can become means of concurrent regulation. Language permits internalized instruction and more precise categorization. Increasing memory and representational capacity lengthen the forecast horizon, allowing action on a later imagined state. Social institutions (families, schools, schedules, alarms, calendars, laws, and built environments) scaffold capacities that no isolated brain exercises alone.

Later forms do not simply replace earlier ones. Adults still seek external soothing, pace, repeat, distract, and discharge. A person capable of long-range planning may revert under acute distress to a shorter regulatory horizon. Such regression does not necessarily erase the acquired capacity; it can indicate that the conditions for its deployment are absent. Conversely, repeated nonuse, disease, injury, or aging can produce more durable decline.

\begin{principle}[Developmental Availability]
Regulatory development is sequential without being irreversible: capacities emerge in an intelligible order, but their availability remains state-dependent and susceptible to regression.
\end{principle}

This avoids two moral errors. A present failure of regulation does not reveal a fixed character, and it does not automatically reveal a single chemical cause. It may reflect which portion of a historically acquired repertoire was accessible in the situation. The relevant causes can include physiology, learned strategies, environmental scaffolds, interpersonal safety, material resources, and the temporal distance at which consequences could be represented.

\section{Surprise, Rehearsal, and the Modeled Possible}

The Principle of Causal Nonidentity applies not only to substances but to the intensity of reactions. A severe reaction may involve endocrine, autonomic, neural, immune, and motor changes. Their involvement does not establish that ``hormones'' explain why this event produced this magnitude of response in this person. The reaction is organized partly by appraisal: whether the event was expected, intelligible, controllable, familiar, and already represented as possible. As in the lottery-ticket case, a real bodily response can be driven by an interpretation whose relation to the external situation remains mistaken or incomplete.

Predictive-processing language offers a useful analogy. An outcome falling far outside an organism's anticipated range generates more surprise, or prediction error, than an outcome already included in its model. The analogy need not commit the essay to a complete predictive-processing theory of mind, nor imply a simple numerical correspondence between surprise and hormonal concentration, and it is worth being explicit about where the commitment actually stops. The claim this essay needs is only that expectation and familiarity modulate the magnitude of a response, which is well supported by the appraisal and expectancy literature already cited throughout; it is a considerably stronger and more contested claim that the brain's fundamental mode of operation is the minimization of prediction error, and nothing that follows depends on that stronger thesis being true. Actual danger, pain, controllability, prior trauma, present health, social support, and learned meaning also shape response. The narrower claim is comparative: familiarity and expectation can help explain why ostensibly similar events affect different people differently, and why the same person may react differently after rehearsal or experience.

A well-studied case shows that single-exposure causal attribution is sometimes exactly the right analysis rather than a just-so story, and the exception is instructive about what the rest of this essay is really objecting to. Conditioned taste aversion, first demonstrated when rats given saccharin during exposure to gamma radiation later avoided saccharin altogether, can be acquired in a single trial, tolerates delays of hours between tasting and illness, and selectively binds a gustatory cue to subsequent gastrointestinal malaise rather than to co-occurring visual or auditory cues encountered at the same time \citep{garcia1955conditioned,garcia1966relation}. A person who becomes unable to tolerate a food after vomiting shortly afterward is not committing the error this essay otherwise criticizes. The mechanism is independently established, its timing and selectivity match its known signature, and a single trial is precisely the condition under which the mechanism is expected to operate. The evidential standard developed above does not forbid strong inference from a single case; it requires that the inference match an independently supported mechanism rather than borrowing the appearance of one.

Three forms of pre-exposure can enlarge what is modeled as possible. The first is \emph{dose titration}: encountering a material in amounts small enough to remain tolerable and increasing exposure as familiarity and competence develop. A food that has become aversive, whether through a conditioned taste aversion or through ordinary dislike, can sometimes be reintroduced by beginning with a fragment far below the amount that provokes rejection and increasing it as tolerance returns, treating the aversion structurally like a phobia subjected to graded exposure rather than as a fixed trait of the food or the person. One author of this essay reports having used this technique repeatedly after childhood illnesses produced food aversions, motivated less by a wish to eliminate discomfort than by a wish not to lose access to food's pleasures, and reports having remained an unusually non-selective eater into adulthood. That self-report is offered as illustration, not evidence: it is a single life without a comparison condition, precisely the kind of claim the Missing Comparison principle above warns against treating as a general finding about how pickiness forms rather than as one plausible mechanism among others. A clinical exposure protocol is more exacting and should not be improvised for allergens, trauma, or conditions in which exposure carries material danger.

The second is \emph{contextual embedding}. An unfamiliar or disliked food may first be combined on the fork with familiar or preferred foods rather than encountered in isolation, distributing textures and flavors across a single bite rather than eating each element as a separate block. Its sensory identity is not erased; it is placed inside an already admissible pattern. As tolerance develops, the familiar support can be reduced. This differs from dose titration because the crucial manipulation is not merely quantity but relational context. The same amount may be experienced differently when embedded in a familiar mixture. Combining ingredients on a plate can serve a further purpose beyond exposure management: for someone assembling a meal as a deliberate recombination of specific remembered dishes, a grandmother's cooking, a particular street food, the mixture is not only a technique for tolerating an unfamiliar element but a way of making a new plate carry several remembered ones at once. That motivation belongs to personal history and culinary culture rather than to the exposure mechanism itself, but it indicates that contextual embedding need not be undertaken for regulatory reasons alone.

\emph{A speculative excursus on vocabulary.} It would be easy to overstate the gap this section is about to name. Positive flavor conditioning is not an understudied domain; flavor-flavor learning, flavor-nutrient learning, mere exposure, and conditioned preference are all established research areas with their own literatures \citep{zellner1983conditioned,rozin1985pavlovian}. The narrower claim worth defending is not that positive conditioning is unstudied but that no name in ordinary usage picks out its one-trial, long-delay form specifically, the way ``taste aversion'' picks out the negative one-trial case with a long cue-outcome delay. \emph{Taste cohesion} is proposed here for that narrower gap, and the term should be read as naming exactly one of three things a looser use of it might otherwise blur together: an increase in liking produced by embedding a disliked element inside an already-familiar sensory whole (the contextual-embedding mechanism discussed above), the durable upgrading of a food's standing after one memorable positive experience with it, or the more autobiographical case in which a dish is deliberately assembled to carry several specific remembered meals at once. This essay uses the term in the second, narrower sense: single-experience upgrading, not gradual conditioning and not the autobiographical case, both of which are already named by contextual embedding and by ordinary recollection respectively. Existing evaluative-conditioning research documents effects built up over repeated pairings rather than the single strong-exposure event that defines taste aversion, and the underlying effects have not always replicated cleanly \citep{zellner1983conditioned,rozin1985pavlovian}. The asymmetry worth defending is therefore specific: one-trial aversion with long cue-outcome delays is well established and biologically prepared, while a symmetrical one-trial positive preference of comparable robustness has not been demonstrated. That asymmetry may be principled rather than a mere gap in study design: one violent illness carries an obvious selective rationale for immediate, one-trial, long-lasting avoidance, since a single misidentification of a toxic food could be fatal, while one good meal carries much weaker pressure toward an equally immediate and permanent preference. Proposing the term is therefore a claim about missing vocabulary for a narrow and real asymmetry, not a claim to have found a hidden mirror image of the Garcia effect, and not a claim that positive flavor learning in general lacks a research literature.

The third is \emph{simulated pre-exposure}. Reading, film, storytelling, play, and deliberate rehearsal can populate a person's model with illness, death, betrayal, violence, disaster, loss, moral conflict, and unfamiliar social worlds before a corresponding event occurs directly. Fiction does not reproduce the event, and imaginative acquaintance must not be equated with having endured it. It can nevertheless make a category of event less wholly unrepresented. The later occurrence may remain new as an episode while no longer being inconceivable as a kind.

This proposal sits near several literatures that should not be collapsed. One treats fiction as a simulation of social experience: narratives abstract relations among agents, intentions, conflicts, emotions, and consequences, allowing readers to run models of social worlds without literally entering them \citep{mar2008fiction,oatley2016fiction}. That is a theory of what narrative engagement can provide, not proof that reading about a catastrophe inoculates a person against experiencing one. A second literature studies deliberate mental simulation as a self-regulatory procedure. Imagining processes, obstacles, or controllable future actions can sometimes improve planning and coping, while fantasy concentrated on a desired endpoint can have different effects \citep{taylor1998imagination,rivkin1999mental}. A third studies stress-exposure or stress-inoculation training, in which people encounter designed stressors while practicing recognition and management of their responses. There the exposure is structured, debriefed, and linked to skills; it is much closer to rehearsal than ordinary reading is. Evidence from such training cannot simply be transferred to fiction, but it demonstrates why controllability, dosage, feedback, and opportunity to practice matter.

The modest claim about fiction is semantic and anticipatory before it is therapeutic. Stories can add cases to the set of situations a person recognizes as humanly possible. They can supply vocabulary for states not yet personally encountered, display several responses to the same predicament, and reveal that an event has an aftermath rather than being an unintelligible terminus. This may reduce one component of surprise: not pain, danger, or grief itself, but the additional disorganization produced when a category is absent from the person's model. The claim is most plausible when narrative acquaintance changes recognition and interpretation. It is much less secure when it is extended to physiological resilience, procedural skill, or protection from trauma.

The distinction between \emph{semantic}, \emph{procedural}, and \emph{affective} preparation is therefore essential. Semantic preparation means knowing that a kind of event exists and having concepts with which to describe it. Procedural preparation means knowing what to do and being able to do it under pressure. Affective preparation would mean that the later emotional response is more tolerable or better regulated. A novel can plausibly furnish the first, occasionally contribute to the second by representing options, and still fail to provide the third. Someone may understand bereavement with great literary subtlety and remain devastated by a death. The devastation does not show that the reading did nothing, and the reading does not entitle anyone to claim equivalence with loss.

Transfer is the hard problem. A simulated situation omits much that organizes a real event: bodily pain, sensory intensity, personal stakes, irreversibility, responsibility, duration, social exposure, and the need to act before reflection is complete. Deliberate rehearsal transfers best when the practiced cues, decisions, and responses resemble those later required. Fiction is typically less constrained. Its power may lie precisely in recombination and ambiguity rather than faithful task simulation. It can broaden the imagination while teaching the wrong procedure, and it can make a scenario vivid without making the reader competent inside it.

Nor is more exposure automatically better. Narrative can sensationalize rare dangers, make availability feel like probability, furnish stereotypes, normalize cruelty, or repeatedly confirm helplessness. Direct exposure can habituate, sensitize, traumatize, or teach avoidance depending on controllability and outcome. Even formal simulation can overload learners if intensity exceeds their capacity to process and recover. Pre-exposure works as regulation only when it expands tolerable prediction and response rather than making threat more chronically available.

The progression also requires terminological precision. A strictly monotonic exposure ladder never decreases in difficulty. If a person returns to an easier stage after distress and later resumes, the overall course may be graded, directional, or ascending in tendency, but it is not mathematically monotonic. Regressions need not be failures; they may be feedback that keeps learning within a tolerable range. This is consistent with the ontogenetic-parade model: development can be ordered without being smooth or irreversible.

Taken together, titration, embedding, and simulation show agency acting before a full reaction takes control, as the person changes dose, context, or anticipated possibility without thereby removing chemistry from the event; what changes instead is the set of conditions under which that chemistry participates. ``My hormones did this'' then becomes an incomplete description of implementation, while ``this event exceeded what I could anticipate and regulate under these conditions'' identifies a more discriminating causal problem. Familiarity is not always the dominant variable, but unlike a retrospective chemical label, it is one variable that can sometimes be deliberately cultivated.

\section{From Possible Hazard to Ordered Risk}

Folk health discourse often asks whether something \emph{can} cause harm. The question is almost designed to produce alarm. At some dose, by some route, in some susceptible organism, under some condition, an enormous range of ordinary materials can cause harm. Possibility identifies a hazard; it does not rank a practical risk.

Risk ordination asks comparative questions: what the magnitude of the effect is at the exposure actually encountered, how probable it is, whether the harm is transient, cumulative, reversible, or catastrophic, which populations are especially susceptible, and what is sacrificed by avoiding the exposure altogether, all weighed against how the risk in question compares with sleep deprivation, alcohol, smoking, infection, sedentary behavior, social isolation, financial precarity, domestic violence, or the risk introduced by the proposed remedy itself.

This comparison is not a license to dismiss small or uncertain harms, since large populations can make small individual effects important, rare catastrophic outcomes can merit precautions on their own, and people may reasonably value risks differently from one another. But a health recommendation that magnifies a trace dietary exposure while ignoring a dominant structural risk has not become rational merely because the trace risk is conceivable.

Commercial attention systematically distorts ordination. Products are easiest to sell when the cause can be placed in a purchasable or avoidable object. ``Control cortisol,'' ``balance hormones,'' ``detoxify,'' ``remove inflammatory foods,'' and ``support neurotransmitters'' convert diffuse problems into consumer tasks. Housing insecurity, exhausting work, caregiving burden, abusive relationships, inaccessible health care, and lack of time are harder to package. Molecular explanations can therefore individualize conditions whose largest causes are infrastructural.

The error is not necessarily conspiracy. A supplement seller, influencer, clinician, and consumer may all sincerely inhabit the same causal vocabulary. Markets select explanations that imply transactions. The resulting advice can contain useful practices (sleep, exercise, adequate nutrition, medical evaluation) while misdescribing why they help or overstating the precision with which a particular molecule has been targeted.

\section{Finger Painting Before Blueprint}

Before presenting the full evidential standard, it is worth explaining why this essay has approached its subject the way it has: introducing the same handful of examples repeatedly, at increasing resolution, rather than presenting the complete technical picture once. This is not incidental style. It follows from a structural fact about why false chemical explanations remain competitive even when a true, more complete explanation is available and well supported.

A false chemical account is almost always cognitively cheap. ``MSG is a neurotransmitter, so it overstimulates your brain'' fits in one breath. A competent correction is not similarly compressible: it has to move through chemical identity, free versus protein-bound fractions, dose, intestinal handling, circulation, barrier transport, extracellular concentration, receptor access, homeostatic regulation, study design, expectancy, and individual variability, most of which this essay has now covered at some length. By the time the correction has become accurate, it can sound evasive, hedging, or needlessly complicated next to the false claim's one confident sentence. This asymmetry is worth naming directly:

\begin{principle}[Explanatory Bandwidth Asymmetry]
A causal chain can be erased in a single sentence but may require an ordered sequence of models to reconstruct.
\end{principle}

Deliberate misinformation sometimes exploits this asymmetry on purpose, overwhelming a listener with more false claims than can be individually addressed in the time available, a tactic sometimes called the Gish gallop after its use in debate contexts. But the asymmetry helps ordinary, non-deliberate misinformation just as effectively, because a compressed falsehood will always look more decisive than an accurate account that has to restore everything the falsehood quietly deleted.

The tempting answer, maximal disclosure delivered all at once, does not solve this problem; it risks reproducing the Gish gallop's functional effect even when every individual sentence is true. A reader confronted with a dozen unfamiliar distinctions simultaneously, compartmentalization, receptor saturation, mediators versus moderators, non-monotonic dose-response, cannot yet tell which distinctions govern which others, and more information under those conditions produces less orientation rather than more. This essay's own length is, on this diagnosis, a real risk rather than a neutral fact about thoroughness: a long document about the danger of collapsing complexity into a single sentence can still fail its reader if it delivers that complexity as an undifferentiated mass rather than as something navigable.

The alternative is to stage complexity rather than either erase it or dump it. Call the first pass at any complicated causal claim a finger painting: imprecise compared to a technical blueprint, but sufficient to establish shape, orientation, and the relations that will matter later. The MSG-lottery pairing that opens this essay is exactly this kind of finger painting:

\begin{quote}
A real chemical exposure does not guarantee the attributed bodily effect, and a false interpretation can produce a real bodily response.
\end{quote}

That single pairing is enough for a first encounter with the whole argument. A next pass separates exposure, bodily response, and causal attribution as three distinct things rather than one fused claim. A further pass introduces the missing physiological gates: digestion, circulation, barrier transport. Still further passes add compartmentalization specifically, then counterfactuals and the absence of a private control condition, then expectancy and self-confirming interpersonal loops, then the causal-role taxonomy, then risk ordination. Each pass returns to material the essay has already touched, at higher resolution, rather than discarding it for something new. A blueprint, on this account, is not a different document from the finger painting; it is the finger painting redrawn with everything the earlier pass could not yet show.

This structure requires its own discipline, because ``the reader now understands more'' is too weak a standard on its own; a reader can accumulate more facts while losing track of which distinctions those facts were supposed to protect. What should not regress, even while specific claims are revised, qualified, or complicated, is the small set of distinctions each pass exists to install:

\begin{principle}[Monotonic Directionality]
An explanation may proceed through simplification, revision, hesitation, and local qualification, but each successive stage should preserve the distinctions already established rather than allowing the reader to collapse them back into the original causal error.
\end{principle}

This is a claim about directionality, not about the total volume of detail a reader currently holds, and the two should not be confused. A reader partway through this essay may not yet know how glutamate transporters work, but if the earlier passes have done their job, that reader already knows that ``glutamate is a neurotransmitter'' does not entail ``eating it changes brain chemistry,'' and no later, more technical pass should require that distinction to be relearned from scratch. That is genuine progress, because it changes which inferences the reader is willing to license, even before every technical detail behind the distinction has been supplied.

This structure is a spiral rather than a ladder. A ladder discards each rung once climbed past; a spiral returns to the same landmarks repeatedly, seeing more of them each time. It is worth distinguishing this pedagogical spiral explicitly from the ontogenetic parade discussed earlier in this essay, because the two are related without being the same thing. The ontogenetic parade describes when a capacity or a fear emerges, plateaus, regresses, and declines over a person's development; it is a claim about developmental timing. Finger Painting Before Blueprint describes how an explanation should be staged for a learner regardless of that learner's developmental stage; it is a claim about pedagogical sequencing. The two interact, since what counts as an admissible next pass depends on which conceptual distinctions and prior models a given reader already has available, but collapsing one into the other would lose what is distinctive about each.

The remaining sections of this essay, the practical evidential standard and the worked examples that follow it, are the blueprint this finger painting has been preparing the reader to receive. They are more technical and more exacting than anything before them, and that is by design: the earlier passes exist so that this final level of detail arrives as a refinement of a shape the reader can already see, rather than as an undifferentiated wall of qualification.

\section{A Practical Standard for Chemical and Hormonal Claims}

The alternative to folk causation cannot be an instruction to wait for perfect certainty. It must be a usable method for grading claims. The following standard is sequential but not mechanical; different questions will require different designs.

\subsection{Specify the outcome}

``Affects hormones,'' ``causes inflammation,'' ``changes brain chemistry,'' and ``disrupts metabolism'' are not adequate endpoints, since nearly any meal affects metabolism and many ordinary experiences alter endocrine activity to some degree. A specified endpoint has to say what changed, by how much, how it was measured, and over what interval, and it has to identify whether the endpoint in question is a laboratory value, a symptom report, a diagnosis, a behavior, or a long-term outcome, because these are different kinds of claims with different evidentiary requirements. Vagueness allows a true triviality to borrow the importance of a much stronger unproven claim.

\subsection{Establish the relevant exposure}

Detection is not dosage. The chemical form, route, concentration, duration, frequency, vehicle, and co-exposures must resemble the claim being evaluated. Injection evidence does not transfer automatically to eating. A fasting bolus does not automatically represent seasoning in a mixed meal. An animal dose many times ordinary human exposure may reveal a hazard while saying little about common risk.

\subsection{Trace a credible route to the target}

The account must pass through the relevant physiology: digestion, transformation, absorption, first-pass metabolism, circulation, binding, clearance, barriers, and transport. This step is where compartmentalization does its critical work. A proposed pathway need not be completely understood before an association is taken seriously, but a known impossibility or grossly implausible exposure cannot be repaired by repeating the molecule's name.

\subsection{Measure perturbation at the proposed target}

Arrival is not action. Material may reach a tissue without reaching relevant cells, may bind without a meaningful response, or may produce a response absorbed by regulation. Biomarkers require validation: a measurable change can be statistically reliable yet biologically trivial, downstream rather than upstream, or unrelated to the outcome that concerns the person.

\subsection{Connect perturbation to experience or conduct}

Even a genuine biochemical change may be one contributor among many, a correlate, or a consequence. Behavioral endpoints are multiply realizable: similar conduct can arise through different paths, and similar biochemical changes can accompany different conduct. The stronger the move from molecule to complex action, the more intermediate evidence is required.

\subsection{Construct a comparison}

What happened under placebo, blinding, withdrawal and rechallenge, a different dose, a matched control, repeated within-person observation, or a natural experiment? The goal is not ritual adherence to one privileged design but approximation of the missing counterfactual. Every design should state which alternative explanations it excludes and which remain.

\subsection{Control the observer and the hypothesis}

Expectation can change symptom reporting, attention, physiology, clinical judgment, and interpersonal conduct. Blinding blocks a real causal path; it is not bureaucratic decoration. Where blinding is impossible, objective measures, independent raters, preregistered outcomes, transparent exclusions, and adversarial tests become more important.

\subsection{Demand stability and convergence}

Demanding stability and convergence means asking whether the result reproduces, whether it persists under reasonable changes in measurement and analysis, and whether mechanistic, observational, and experimental evidence converge on the same conclusion, together with whether the relation follows a coherent dose-response pattern and whether null results are available alongside the positive ones. A single spectacular finding should normally receive less weight than a modest relation that holds up stably across several independent methods.

\subsection{Ordinate the practical risk}

Finally, how does the estimated magnitude and probability compare with competing risks, opportunity costs, and intervention burdens? A real but tiny effect can be a poor target. A low-cost personal precaution may be rational under weak evidence, while a restrictive public recommendation demands stronger support. Evidential thresholds should vary with the costs of acting and failing to act.

This standard does not issue automatic verdicts. It disciplines the route by which verdicts are reached. Most importantly, it applies symmetrically. A reassuring claim must establish its scope just as an alarming claim must establish its mechanism. ``Safe'' should mean safe under specified conditions and evidence, not metaphysically incapable of harm.

The standard also separates two questions that folk discourse tends to run together: evidential sufficiency and decision sufficiency. Evidential sufficiency asks whether the nine questions above have been answered well enough to support a general causal claim. Decision sufficiency asks whether the available evidence, however incomplete, is enough to justify a particular action given its cost. A person may lack evidence sufficient to claim that MSG generally causes headache in the population, while possessing personal evidence sufficient to avoid it at negligible cost to themselves. A public-health prohibition or a clinical treatment recommendation requires a different, generally higher threshold, because its burdens fall on people who have not personally observed the pattern and because reversing a population-level policy is more costly than reversing a private habit. Confusing these two kinds of sufficiency runs in both directions: it produces both the person who treats a personal pattern as license to make sweeping claims to others, and the critic who withholds a cheap personal precaution because it has not cleared the higher bar appropriate to policy.

\subsection{Three claims run through the standard}

Running distinct claims through all nine questions shows that the standard is a portable method rather than a rule invented for MSG alone.

\emph{``MSG causes headaches.''} The outcome is specifiable (a headache, timed and rated). The exposure is not homogeneous: culinary MSG dispersed through a meal is a materially different exposure from a large aqueous bolus consumed while fasting, and most existing challenge studies used something closer to the latter. A route to the target is plausible for some symptoms without requiring central nervous system penetration at all, since taste, expectation, and peripheral effects are available candidate mechanisms; a claim requiring MSG glutamate to cross into brain tissue in physiologically significant amounts faces the barrier evidence reviewed above. Perturbation and its connection to the reported symptom have been measured only inconsistently, with the most rigorous multicenter trial finding limited reproducibility on rechallenge \citep{geha2000multicenter}, and a systematic review of headache specifically finding pervasive blinding failures caused by MSG's detectable taste \citep{obayashi2016headache}. A comparison condition (placebo, blinding) exists in the better studies and is precisely where most of them fail. Expectation is a serious threat to validity here, given MSG's cultural salience. Stability across studies is weak. Risk ordination favors a narrow personal-avoidance policy for someone who reliably notices a pattern with large culinary doses, and disfavors a general population-level warning about ordinary culinary use.

\emph{``Hormones in chicken alter human behavior.''} The outcome (some behavioral change) is rarely specified beyond vague reference. The exposure conflates the three propositions distinguished above: endogenous hormones are certainly present, but added hormonal growth promoters are not permitted in poultry in jurisdictions such as Canada in the first place, so the exposure the claim implicitly assumes may not exist as described. Even granting endogenous hormone presence, a route to a human behavioral target requires the hormone to survive cooking and digestion, reach the bloodstream in an active form, and do so in an amount large relative to the consumer's own endogenous production, a chain with no public evidence offered in most versions of this claim. No perturbation at a human target has generally been measured. No comparison condition is typically constructed. The claim fails at nearly every stage before risk ordination becomes relevant at all, which is itself diagnostic: a claim that cannot be evaluated because the necessary studies were never done is different from a claim that has been tested and found weak, and the practical response (seek the relevant exposure data, or withhold judgment) differs accordingly.

\emph{``My cortisol caused this decision.''} The outcome, a decision, is specified, but ``my cortisol'' as a cause is not, since it fails to specify basal, reactive, or chronic profile, or by what measure it was assessed, as the cortisol subsection above discusses. No exposure or dose question applies in the usual sense, since this is an endogenous-state claim rather than an ingestion claim; the corresponding question becomes whether the person's cortisol was actually measured at the relevant moment, which in nearly all everyday invocations it was not. A route from an elevated hormonal state to a specific complex decision is not established merely by cortisol's general association with stress physiology; decisions are downstream of appraisal, available options, social context, and habit as much as of endocrine state. No comparison condition exists (the same decision moment cannot be rerun with cortisol experimentally lowered). The claim functions more as a retrospective narrative, in the sense developed earlier, than as a tested causal proposition, and evaluating it well means asking what alternative account of the decision, involving triggers, permissive conditions, and prior commitments, would need to be ruled out before the hormonal narrative could be preferred over them.

The three examples do not fit the standard equally well, and that unevenness is itself worth naming rather than treating as an accident of selection. The MSG claim is a population-level empirical proposition with an actual, if imperfect, body of trials behind it, so the nine questions above can be run against real data at every stage. The chicken-hormone claim fails earlier, often at the exposure stage itself, because the relevant studies were largely never conducted in the first place; that is a different evidential situation from a claim that was tested and found weak, and it calls for seeking the missing exposure data rather than weighing existing results. The cortisol claim is not a population-level proposition at all but a first-person retrospective narrative about a single decision, which is exactly the kind of claim the Missing Comparison principle earlier in this essay says cannot be evaluated by the same methods used for the other two. A single evidential standard applied honestly will therefore fit these three claims differently, and recognizing which kind of claim is on the table, tested-and-weak, untested, or retrospective-and-uncomparable, is as much a part of using the standard well as running through its individual questions.

\section{Chemical Stories as Cognitive Strategies}

Why do weak chemical explanations persist after correction? One answer is informational: the science is difficult, public communication is poor, and memorable simplifications spread readily. That answer is incomplete because it treats the explanation only as a failed representation of physiology. Often it is also a successful cognitive strategy.

``MSG caused my headache'' takes a diffuse, underdetermined event and gives it a boundary. ``My hormones made me do it'' divides the person into an enduring self and a temporary chemical force. ``Cortisol is causing my belly'' gathers appearance, anxiety, aging, stress, and social pressure into a named object. Each formulation reduces causal entropy. It identifies a victim, an agent, and an intervention: eliminate the ingredient, correct the hormone, buy the supplement, change the routine, regain control.

Medical anthropology and health psychology offer more precise ways to understand this work. Kleinman, Eisenberg, and Good used the term \emph{explanatory model} for the understandings of illness held by patients, families, and practitioners: accounts of what the problem is, why it began, how it operates, how severe it is, what course it will take, and what should be done \citep{kleinman1978culture}. Such models are not miniature scientific papers stored in people's heads. They are practical orientations assembled from bodily experience, family knowledge, clinical language, media, religion, markets, and local moral expectations. They guide what counts as a symptom, when help is warranted, who is responsible, and which remedy is sensible.

The common-sense model of illness self-regulation describes a related process in more explicitly psychological terms. People form cognitive representations of a health threat along dimensions such as its \emph{identity}, proposed \emph{cause}, expected \emph{timeline}, anticipated \emph{consequences}, coherence, and perceived possibilities for control or cure. Emotional representations develop alongside them. These representations guide coping procedures; their outcomes are appraised; the model and the coping response can then be revised in a continuing loop \citep{leventhal1980common,hagger2017commonsense,hagger2022conceptual}. A large meta-analysis found patterned relations among representation dimensions, coping strategies, distress, functioning, and disease outcomes, while also showing that the model is not reducible to a single linear belief-behavior rule \citep{hagger2017commonsense}.

Chemical stories fit this architecture remarkably well. ``Cortisol dysregulation'' supplies an identity and a mechanism-like cause. ``Perimenopause'' or ``inflammation'' supplies a timeline, sometimes a beginning and anticipated endpoint. Fatigue, weight change, irritability, pain, fogginess, and relationship strain become consequences of one named process. Supplements, elimination diets, exercise prescriptions, laboratory panels, or hormone treatment become control procedures. Online testimony provides coherence by teaching which experiences belong together. The story is attractive not only because a molecule sounds scientific but because it fills nearly every practical slot an illness representation requires.

This does not make all such models false. An explanatory model may accurately identify a disease, organize effective treatment, and articulate an experience that medicine previously neglected. Perimenopause, endometriosis, postpartum depression, and adverse drug effects have all been vulnerable to dismissal when testimony did not fit established clinical attention. The argument here is not that lay interpretation is intrinsically inferior to professional interpretation. Clinicians also use incomplete models, and patients may detect patterns that short appointments miss. The point is that usefulness, professional status, vividness, and mechanistic accuracy are separable properties.

Naming can be genuinely beneficial. A diagnosis can secure recognition, organize care, connect people with similar experiences, justify accommodation, and replace shame with intelligibility. Externalization can create the distance needed for action: the person is not simply lazy, weak, irrational, or morally defective. Yet the same operation can over-externalize agency, turn a contributing condition into a sovereign cause, or make every discordant event evidence of the category. The mistake would be to infer that psychological usefulness validates the proposed biological mechanism. A category can regulate emotion while misrepresenting causation.

An explanatory model also performs \emph{compression}. A person cannot ordinarily represent digestion, receptor trafficking, labor conditions, childhood learning, sleep debt, household conflict, financial pressure, and cultural ideals every time a headache or argument occurs. A chemical noun drastically reduces the number of variables under active consideration. Compression is unavoidable; action would be impossible if every decision required a complete causal graph, the ordinary version of the bandwidth problem discussed in the preceding section. The relevant question is what the compression preserves and what it discards. ``MSG'' may preserve temporal association and a feasible avoidance strategy while discarding meal composition, expectancy, headache base rates, and dose. ``Hormones'' may preserve the reality of cyclical or transitional bodily variation while discarding a partner's conduct, economic dependence, training, and institutional support.

The model can further distribute moral responsibility. ``My hormones made me do it'' may protect a valued self from identification with an act. ``Her hormones made her do it'' may instead invalidate her reasons and license another person's refusal to listen. The same grammar therefore can be compassionate in first-person use and coercive in third-person use. It can reduce self-blame, intensify stigma, excuse harm, or deny grievance depending on who controls the attribution. Causal vocabulary is never only descriptive once it enters a relationship.

Chemical stories can also protect relationships and institutions. Blaming sugar may be easier than examining a chaotic party, inconsistent boundaries, exhaustion, or conflict. Blaming hormones may be easier than hearing a partner's grievance. Blaming individual cortisol may be easier than changing a workplace. A tractable molecule can become the designated cause precisely because larger causes are socially expensive. When the intervention is privatized, purchase the supplement, optimize the meal, discipline the body, the explanatory model can preserve the conditions that repeatedly generate distress.

Markets do not merely attach products to pre-existing beliefs. They help stabilize the representations through which symptoms are noticed. A feed filled with ``cortisol face,'' ``menopause belly,'' glucose spikes, detoxification, and hormone-balancing routines teaches an identity, a symptom inventory, a causal account, and a control procedure at once. The buyer then monitors experience using the supplied categories. Some will obtain useful recognition or adopt beneficial routines. Others will acquire anxiety, unnecessary restriction, false reassurance, or a delay in care. The commercial form is powerful because it offers interpretation and action in the same package.

This is why correction by ridicule fails, and it suggests a \emph{replacement burden} that any adequate correction has to discharge. If an explanation supplies orientation, removing it without replacement restores the original disorder. The person loses not only a proposition but a vocabulary, a social community, a division of responsibility, and a plan. ``There is no good evidence that MSG caused this'' may be accurate and still pragmatically incomplete, because it leaves unanswered what the person should notice instead, what can be tried safely, which clinician will take the experience seriously, what uncertainty must simply be tolerated, and which larger condition is actually alterable. A correction that answers none of these questions asks someone to exchange an actionable error for unaided ambiguity.

Respectful inquiry therefore begins by eliciting the model before attacking it, asking what the person calls the problem, what they think caused it, why they believe it began when it did, what they think it does inside the body, what course they expect it to take, what frightens them about it, and what they have already tried and with what result. These questions do not endorse the answers they receive. They reveal the functions that a replacement account must perform and the observations on which the current one rests, and they can also expose mismatches between patient and clinician: one may be trying to explain pain and preserve dignity while the other is trying to classify pathology.

Correction can then separate several propositions that a single label had fused. The experience may be real while the proposed molecule remains unmeasured. A biological transition may matter while not dictating conduct. A personal avoidance may be reasonable while a universal toxicological claim remains unsupported. A treatment may help through a pathway different from the one advertised. A diagnosis may offer belonging while still being too broad for a specific causal inference. This decomposition preserves testimony without granting every explanation attached to it.

The replacement may be practical rather than metaphysically complete. Keep a blinded symptom diary where possible. Change one exposure repeatedly rather than many at once. Record dose, timing, sleep, context, and competing exposures. Distinguish a private avoidance policy from a universal claim. Seek measurements when they would alter treatment. Compare the suspected risk with dominant risks and the remedy's costs. Improve sleep, safety, nutrition, workload, and care access without pretending each intervention precisely ``balances'' a named chemical. The aim is not to eliminate uncertainty before acting but to prevent action from hardening uncertainty into invented mechanism.

The final task is not to make every person speak in provisional causal diagrams. Ordinary life requires names, compression, and action. It is to cultivate models that retain revisability: stories able to say what evidence would count against them, which claims are personal policies rather than universal facts, and where social or material causes enter alongside biological ones. A good explanatory model is not necessarily complete. It is corrigible, proportionate, and generous enough to keep more than one causal level in view.

\section{Conclusion: An Orientation Without a Verdict}

The central error examined here is simple to state and difficult to stop committing: chemical identity substitutes for causal context. The same mistake declares artificial substances dangerous because they are artificial and endogenous substances harmless because they are endogenous. It moves from dietary glutamate to synaptic excitation, from the presence of hormones in a food to altered human behavior, and from a felt state to an unmeasured endocrine diagnosis.

Living systems make these shortcuts unreliable because they are organized by gates. Dose, form, route, metabolism, compartment, transport, receptors, timing, regulation, development, expectation, and social interaction determine which continuations are physically available. The blood--brain barrier is one especially vivid gate, but the principle is general. A molecule's name does not grant passage.

Human experience adds a second limitation. A lived sequence has no private control condition. Introspection can reveal suffering and temporal order without isolating an invisible biochemical cause. Causal methods approximate the missing comparison, and expectancy must be treated as part of the system rather than dismissed as unreality. The lottery ticket shows an interpretation altering the body without the attributed event being real. The sugar-expectancy experiment shows a chemical belief altering observation and interaction even when exposure is held constant.

Agency does not float above these processes. It develops through them. Prospective regulation acts upon a modeled later state; concurrent regulation changes an active state through feedback. Their capacities emerge in a developmentally ordered but non-monotonic succession, remain dependent upon bodily and social conditions, and can regress under strain. A hormone can alter the difficulty landscape without choosing the action. A cognitive strategy can reorganize the landscape without abolishing chemistry.

Practical judgment must therefore replace categorical fear with evidential tracing and risk ordination: asking what the endpoint actually is, whether the exposure was relevant, whether the proposed agent could plausibly reach its target, whether the target was in fact perturbed, and whether that perturbation explains the outcome at all, alongside what comparison approximates the absent counterfactual, whether expectation or surveillance might have helped create the very evidence being cited, whether the result converges across independent methods, and finally how the risk in question ranks among the ordinary risks of an actual life.

Yet the final correction is psychological as much as scientific. Folk chemical attribution often names diffuse suffering, externalizes it, and supplies a feasible action. Its persistence cannot be explained solely by chemical ignorance. It is an instrument of emotional organization and sometimes of interpersonal or institutional avoidance.

A chemical story can therefore be psychologically functional, socially consequential, and biologically false at the same time. Correcting it requires more than winning an argument about a molecule. It requires an orientation capable of acknowledging the body without making chemistry destiny; preserving testimony without mistaking sequence for causation; enabling action without manufacturing certainty; and locating the person not beneath a single chemical command, but within the pathways, histories, relationships, and strategies through which a life becomes causally intelligible.

One further question this essay has raised without answering deserves to be named rather than left implicit. Everything offered here, the evidential standard, the causal-role taxonomy, the practice of eliciting an explanatory model before correcting it, operates at the level of an individual claim or an individual conversation. But this essay's own analysis suggests that chemical stories are sustained at least as much by media, markets, and social conditions as by any single person's reasoning, which means that correcting the causal proposition without addressing the conditions that make the story useful may accomplish comparatively little at scale, in something like the way treating a fever without treating the underlying infection brings only temporary relief. What an adequate public epistemology for this problem would look like, as opposed to a private discipline for evaluating one claim at a time, is a genuine further question this essay has not attempted to answer.

\bibliographystyle{unsrtnat}
\bibliography{chemical-identity-is-not-causation}

\end{document}
